% !TEX program = xelatex
% 电工电子技术 · 第一章第三讲（Chapter One -3）· 精讲讲义（lecture-deepdive 规范）
\documentclass[11pt]{ctexart}
%=====================================================================
% theme-violet.tex —— 讲义精讲 · 紫花主题
% 用法：\documentclass[11pt]{ctexart}
%       \input{assets/theme-violet}
%       \begin{document} ... \end{document}
% 只用 ctexart 自带能力 + xcolor/enumitem/fancyhdr/pifont/amsmath，
% 不依赖 tikz/tcolorbox，保证 MiKTeX 可直接编译。
%=====================================================================

\usepackage[a4paper,left=2.0cm,right=2.0cm,top=2.1cm,bottom=2.3cm,%
            headsep=12pt,footskip=22pt]{geometry}
\usepackage{amsmath,amssymb}
\usepackage{xcolor}
\usepackage{array}
\usepackage{enumitem}
\usepackage{pifont}
\usepackage{fancyhdr}
\usepackage[hidelinks]{hyperref}

%---------------------------------------------------------------- 配色
\definecolor{violet}{RGB}{107,78,143}      % 主色：紫
\definecolor{violetdeep}{RGB}{63,44,92}    % 深紫：标题/强调
\definecolor{petal}{RGB}{183,159,208}      % 花瓣：装饰线、花饰
\definecolor{bandviolet}{RGB}{247,243,251} % 浅紫底：档案/提示框
\definecolor{bandcream}{RGB}{252,249,242}  % 米白底：易错/注意框
\definecolor{rulepurple}{RGB}{217,207,230} % 分隔线
\definecolor{inkgray}{RGB}{92,96,110}      % 次要文字

%---------------------------------------------------------------- 字体
% 有 Noto CJK 用 Noto，没有就退回 ctex 默认（Windows 下为宋体/黑体）
\IfFontExistsTF{Noto Serif CJK SC}{\setCJKmainfont{Noto Serif CJK SC}}{}
\IfFontExistsTF{Noto Sans CJK SC}{\setCJKsansfont{Noto Sans CJK SC}}{}

%---------------------------------------------------------------- 版式
\setlist{nosep,leftmargin=1.9em}
\linespread{1.28}
\setlength{\parskip}{3.5pt}
\setlength{\parindent}{0pt}
\setlength{\tabcolsep}{5pt}
\renewcommand{\arraystretch}{1.28}

%---------------------------------------------------------------- 花饰
% \ding{95} = 五瓣小花，\ding{96} = 星形花；都来自 pifont，必定存在
\newcommand{\orn}{\textcolor{petal}{\ding{95}}}          % 行内小标
\newcommand{\ornstar}{\textcolor{petal}{\ding{96}}}      % 星形花
\newcommand{\flowerrule}{%                                % 居中花饰分隔线
  \par\vspace{7pt}%
  \noindent\makebox[\linewidth]{\color{petal}%
    \rule{0.20\linewidth}{0.6pt}\hspace{8pt}{\small\ding{95}}\hspace{4pt}%
    \raisebox{1.1pt}{\tiny\ding{95}}\hspace{4pt}{\small\ding{95}}\hspace{8pt}%
    \rule{0.20\linewidth}{0.6pt}}%
  \par\vspace{7pt}}

%---------------------------------------------------------------- 标题与节头
\newcommand{\doctitle}[3]{%   #1 主标题  #2 副标题  #3 元信息（可用 \\ 分行）
  \begin{center}
    {\color{petal}\small \orn\enspace #2\enspace \ornstar}\\[7pt]
    {\LARGE\bfseries\color{violetdeep} #1}\\[10pt]
    {\color{petal}\rule{0.55\linewidth}{1pt}}\\[4pt]
    {\small\color{inkgray} #3}
  \end{center}\vspace{2pt}}

\newcommand{\hdr}[2]{%   #1 章节号（如 §0）  #2 标题
  \par\vspace{12pt}%
  \noindent\colorbox{violetdeep}{\color{white}\bfseries\small\;#1\;}\hspace{7pt}%
  {\color{violetdeep}\large\bfseries #2}\par\nobreak
  \vspace{3pt}{\color{petal}\hrule height 1.2pt}\vspace{6pt}}

\newcommand{\subhdr}[1]{%
  \par\vspace{6pt}%
  \noindent{\color{violet}\bfseries #1}\par\vspace{1pt}}

%---------------------------------------------------------------- 提示框
\newcommand{\callout}[1]{% 浅紫底：讲义档案、读法说明
  \par\medskip\noindent
  \fcolorbox{rulepurple}{bandviolet}{\parbox{\dimexpr\linewidth-2\fboxsep-2\fboxrule}{\vspace{2pt}#1\vspace{2pt}}}%
  \par\medskip}

\newcommand{\warnbox}[1]{% 米白底：易错、边界条件
  \par\medskip\noindent
  \fcolorbox{rulepurple}{bandcream}{\parbox{\dimexpr\linewidth-2\fboxsep-2\fboxrule}{\vspace{2pt}#1\vspace{2pt}}}%
  \par\medskip}

\newcommand{\gapitem}[2]{% #1 缺口类型（紫） #2 内容
  \par\vspace{2pt}\noindent
  \hangindent=1.4em\hangafter=1
  {\color{petal}\small\ding{96}}\enspace{\bfseries\color{violetdeep}#1}\enspace #2\par}

% 前提旁注：只放在「需要理解的那一步」后面，给陈述 + 本题代入，不重证。
\newcommand{\aside}[1]{%
  \par\vspace{1pt}\noindent
  \hangindent=1.6em\hangafter=1
  {\small\color{inkgray}%
    \hspace*{0.6em}{\color{petal}$\triangleright$}\hspace{2pt}#1}%
  \par\vspace{2pt}}

%---------------------------------------------------------------- 表格
% 本机 MiKTeX 的 colortbl 与 array 版本不匹配（\rowcolor 会报错），
% 故表格统一用「上粗紫线 + 表头 + 中细线」的三线式，不依赖 colortbl。
\newcommand{\tabtop}{\par\vspace{3pt}{\color{violet}\hrule height 1.2pt}\vspace{1pt}}
\newcommand{\tabmid}{\par\vspace{1pt}{\color{petal}\hrule height 0.5pt}\vspace{1pt}}
\newcommand{\tabbot}{\par\vspace{3pt}{\color{rulepurple}\hrule height 0.5pt}\vspace{8pt}}
\newcommand{\thead}[1]{{\bfseries\color{violetdeep}#1}}

%---------------------------------------------------------------- 页眉页脚
\pagestyle{fancy}
\fancyhf{}
\renewcommand{\headrulewidth}{0.4pt}
\renewcommand{\headrule}{\hbox to\headwidth{\color{rulepurple}\leaders\hrule height \headrulewidth\hfill}}
\fancyhead[L]{\small\color{inkgray}\songti \doctag}
\fancyhead[R]{\small\color{inkgray}\songti \coursetag}
\fancyfoot[C]{\small\color{inkgray}\orn\enspace{\color{violet}\thepage}\enspace\orn}

% 每篇讲义在导言区设这两个值
\newcommand{\doctag}{}
\newcommand{\coursetag}{}
\newcommand{\settags}[2]{\renewcommand{\doctag}{#1}\renewcommand{\coursetag}{#2}}

%---------------------------------------------------------------- 数学记号
\newcommand{\vect}[1]{\boldsymbol{#1}}
\newcommand{\dd}{\,\mathrm{d}}
\newcommand{\ddx}{\frac{\mathrm{d}}{\mathrm{d}x}}
\addtolength{\jot}{2pt}



\settags{C1-3 · 精讲}{电工电子技术}

\begin{document}

\doctitle{1.7 基尔霍夫定律与 1.8 支路电流法}
{两句「代数和为零」怎么落到符号上、独立方程数从哪来、支路电流法的四步怎么走}
{依据：Chapter One -3-2026中白上传.pdf（p1--p67，共 67 页；本讲义页码一律指本册 PDF 物理页，
引用他册时写成 [C1-2 p30] 的形式）\ \textperiodcentered\ 2026-10-23}

%=============================================================
\hdr{§0}{课件坐标、记号与四处拼接}

\subhdr{0.1 这一册讲了什么，各在哪几页}

这一册是三段合成：\textbf{先把 1.6 的电源等效变换收个尾}（p1--p3 的复习题、p4--p5 的人物页），
\textbf{再用四页把「节点、支路、回路、网孔」数清楚}（p6--p10），
\textbf{然后依次讲 KCL}（p11--p17）、\textbf{KVL}（p18--p42）、\textbf{支路电流法}（p43--p56），
最后 \textbf{提前开了 1.9 叠加定理的头}（p57--p67）。

\par\vspace{2pt}
\noindent\begin{tabular}{@{}p{3.1cm}p{1.9cm}p{6.0cm}p{4.6cm}@{}}
\tabtop
\thead{课件小节} & \thead{页码} & \thead{内容} & \thead{本讲义对应}
\\
\tabmid
Review+Preview & p1、p43 & 同一张目录页（1.6 复习 + 1.7/1.8/1.9/1.10 预告） & §0.2
\\
1.6 复习题 & p2--p3 & 两道「Find I by using source transformation」，\textbf{只出题} & §0.2、§7.5
\\
人物页 & p4--p5 & Kirchhoff 与 Bunsen 的生平、光谱分析 & 不入正文（见 §0.2 第（4）条）
\\
Terms & p6--p10 & 四个词的定义 + 四页计数练习（p7--p10） & §1.1--§1.3
\\
1.7.1 KCL & p11--p17 & 提要五项、陈述、例、广义节点（closed boundary） & §2
\\
1.7.2 KVL & p18--p27 & 对象（Loop/Mesh）、绕行方向、陈述、E1、E2 & §3--§4
\\
1.7.2 KVL & p28--p35 & Imaginary loop、端口电压 $U=U_S+RI$、电流反向 $U=U_S-RI$ & §5--§6
\\
1.7.2 KVL & p36 & Find $U_1,U_2,U_3,U_4$（四小题挤一页） & §6.3
\\
1.7.2 KVL & p37--p42 & Find $U_{AB}$ 一页、Find $U_{ab}$ 六页（\textbf{全无答案}） & §5.4、§6.4
\\
1.8.1 引 & p44--p47 & 目的、工具（KCL+KVL）、方法（对节点/对回路） & §7.1
\\
1.8.2 步骤 & p48 & 四步（定方向 $\to$ 节点 KCL $\to$ 回路 KVL $\to$ 求解） & §7.2
\\
1.8.3 例题 & p49--p54 & 五支路两电源例题，走到第（4）步就停 & §8.1--§8.3
\\
1.8.3 例题 & p55--p56 & 两张数值练习图（40/20/4 与 5V/10V 六电流） & §8.5、§8.6
\\
1.9 叠加定理 & p57--p67 & 分节页、定义、内容、短路/开路、步骤、注意 & §9
\tabbot
\end{tabular}

\subhdr{0.2 四处拼接：课件自己的缝，与范围粗测的出入}

\gapitem{（1）粗测地图与实物页码不符}{任务给的粗测是「1.7 KVL 及推广 p26--p50 前后、1.8 支路电流法 p50 之后至 p67」。
逐页核对后的实物是：\textbf{KVL 从 p18 开始}（p18 的页眉已经是 1.7.2），
\textbf{1.8 从 p43/p44 开始}（p43 是复习/预习目录页，p44 才是 1.8.1 的第一页），
\textbf{p57 起已经是 1.9 叠加定理}——不是 1.8。本讲义一律按实物页码引用，§12 有对照表。}

\gapitem{（2）四处版面重复}{p1 与 p43 是同一张目录页；p3 与 p45 是同一张电路图
（1.6 的复习题与 1.8 的入门例题用的是\textbf{同一个网络}：$I_S,R_1,R_2,R_3,R_4,U_S$）；
p45 与 p46 只差「电流参考方向箭头」；p61 与 p64 只差「电源符号画回去」。
§0.1 的坐标表按去重后的内容计，§10 记两条。}

\gapitem{（3）两套符号练习之间没有桥}{p21--p23 练「给四个电阻标极性、判正负」，p28--p36 练「端口电压与电流方向」，
两边都只用图说话，\textbf{中间没有一句话把规则写出来}。这是本册最大的机制缺口，§3.2 补一条统一规则，
并\textbf{在 p23、p27、p32、p35 四处独立验证}它——四处都对，才敢说它是规则而不是巧合。}

\gapitem{（4）六处只出题不给答案}{p2/p3（1.6 复习题）、p36（$U_1\sim U_4$）、p37--p42（$U_{AB}$、$U_{ab}$）、
p54（第（4）步「求解」）、p56（六电流图）。本讲义把这六处逐一补上或说明为什么补不了：
p37 与 p38--p42 给出数值（§5.4、§6.4），p36 给出对应表（§6.3），p54 给出矩阵与唯一性（§8.3），
p56 标 [待确认] 并只做方程数讨论（§8.6）；p2/p3 因原图接线在文字层上无法唯一确定，也标 [待确认]（§7.4）。
p4--p5 是 Kirchhoff 与 Bunsen 的人物与光谱分析页，与本册的定律用法无关，
本讲义只在 §0.1 的坐标表里给它们留一行，不再复述生平。}

\callout{\textbf{本讲义怎么用}：\\
第一次读只看 §0.1 的页码表、§8.2 的数值全解、§9.5 的公式卡——这三处能让你把题做对；\\
第二次读再看 §3.2（统一符号规则）、§7.3（独立方程数）、§8.3（方程组为什么有唯一解）——
这三处解释「换一个电路还对不对」，是考试里真正的分界；\\
\textbf{§10 缺口清单}逐条对应课件某页某句：课件哪句话让你卡住，就去那里找它在哪一节补的。}

\subhdr{0.3 记号：先统一，后文不再回头}

\par\vspace{1pt}{\small
\noindent\begin{tabular}{@{}p{2.5cm}p{12.4cm}@{}}
\tabtop
\thead{记号} & \thead{约定}
\\
\tabmid
$i,\ u$ & \textbf{瞬时值}（小写）：只在本册 p13、p20 的陈述句「$\Sigma i=0$、$\Sigma u=0$」里出现，
意思与下一条的直流量相同，本讲义后文一律用大写
\\
$I,\ U$ & \textbf{直流量的参考值}：$I_k$ 是第 $k$ 条支路电流的\emph{代数值}，
它的正负由「参考方向」决定，不代表实际方向
\\
箭头 $\uparrow\downarrow\leftarrow\rightarrow$ & \textbf{参考方向}：图上画的箭头是解题者自己定的，
算出负号只说明实际方向与箭头相反，不说明算错了
\\
$\to$ 与 $\circlearrowright$ & 前者是元件旁的\textbf{元件箭头}，后者是回路里的\textbf{绕行方向}（p19 的 Loop direction）
\\
$U_S,\ I_S$ & 理想电压源、理想电流源的\textbf{源值}，本身恒正；极性由旁边的 $+/-$ 记号给出
\\
$R_k$ & 第 $k$ 条支路的电阻。一条支路可以同时含电阻与电源（如 p26 的 $U_{S1}+R_1$）
\\
$N,\ B,\ M,\ L$ & 节点数、支路数、网孔数、回路数（p7--p10 的 N、B、M、L）
\\
$U_{AB}$ & 端口电压，$U_{AB}=V_A-V_B$；参考方向由箭头给出（p37 的箭头由 A 指向 B）
\\
$[\,$待确认$\,]$ & 原图分辨率下无法唯一确定的事实，不猜、不代填
\tabbot
\end{tabular}}

%=============================================================
\hdr{§1}{四个词与四个数：$N$、$B$、$M$、$L$}

\subhdr{1.1 p6 的四个定义，以及每个定义里那句容易被跳过的话}

课件 p6 一页只有四行，每行一个词加一句解释。这四行是本册全部方程的基础，
所以逐句读，并把「容易跳过的那半句」挑出来：

\par\vspace{2pt}
\noindent\begin{tabular}{@{}p{1.9cm}p{5.3cm}p{8.4cm}@{}}
\tabtop
\thead{词} & \thead{课件原文（p6）} & \thead{本讲义的读法，以及容易跳过的那半句}
\\
\tabmid
Node & Point where three or more elements meet. & \textbf{三条及以上}元件的交点才算节点。
两个元件首尾相接的那一点\emph{不是}节点——这条限定直接决定 $N$ 怎么数，
p7 的三支路并联电路因此只有 2 个节点而不是 4 个
\\
Branch & Circuit branch between two nodes & 两个节点之间的一段，\textbf{一条支路只有一个电流}。
这是「支路电流法」这个名字的来源：未知数就是支路数 $B$
\\
Loop & Any closed path in a circuit & 任意闭合路径，\textbf{不重复走同一段}。
注意：回路不是「网孔」，一个电路里回路的条数 $L$ 远多于网孔数 $M$（见 §1.3）
\\
Mesh & Independent loop. & 这句是四个定义里最省的一句：网孔是「独立回路」，
但\textbf{课件没有说「独立」到底独立在哪}、也没有说网孔有几个。
§1.2 把这两件事补上
\tabbot
\end{tabular}

\aside{p6 的原文用的是「three or more elements」，而 p7--p10 的答案里
$N$ 数的是「三条及以上支路的交点」。两种说法在二端元件电路里等价；
若电路里有三端器件（如后面的三极管），课件的措辞要按「三条及以上支路」来读，
否则会多数。本册不涉及三端器件，按「支路」数即可。}

\subhdr{1.2 四个数之间的关系：本册最重要的一条定理}

课件把 $N,B,M,L$ 的答案直接印在 p7--p10 上（如 p7 的「N: 2\quad B: 3\quad L: 3\quad M: 2」），
但从头到尾没有一个字说明这些数之间有什么约束。补上它，因为 §7 的「独立方程数」全靠它。

\warnbox{\textbf{定理（独立方程数）}\quad 设电路连通，有 $N$ 个节点、$B$ 条支路，且每条支路只有一个未知电流。则：
\begin{enumerate}
\item KCL 能给出的\textbf{独立}方程恰好 $N-1$ 个（任取 $N-1$ 个节点即可，去掉的那一个方程可由其余方程相加得到）；
\item 独立回路（网孔）的个数为
\[ M=B-(N-1)=B-N+1 ; \]
\item 两者合计 $(N-1)+M=B$ 个独立方程，\textbf{恰好等于未知支路电流的个数}——所以方程组有唯一解。
\end{enumerate}}

\subhdr{这条定理的原形式与课件的用法}

要按硬性要求说清「课件给的是哪一个形式」：\textbf{原形式是「KCL 的独立方程数 $N-1$、KVL 的独立方程数 $=B-(N-1)$」}
——它讲的是「独立」二字，不是「$B-N+1$」这个算式。
课件的 $M=B-N+1$ 是把它代进去后的\textbf{常用形式}（在平面电路里，这 $M$ 个独立回路正好就是网孔；一般连通电路里它仍是独立回路的个数，见 §1.6）。
课件从头到尾只用常用形式（p7--p10、p49 的「5 equations」、p52 的「Only two of them are independent」），
原形式一次也没写。

\aside{为什么第 1 条成立（不重证，要说清机制）：设 $\vect{A}_a$ 是 $N\times B$ 的节点--支路关联矩阵，
它的每一行求和为零，所以 $\operatorname{rank}\vect{A}_a\le N-1$；又因为电路连通，
任取 $N-1$ 行都线性无关，所以 $\operatorname{rank}\vect{A}_a=N-1$。
KCL 就是解 $\vect{A}_a\vect{I}=\vect{0}$，解空间的维数是 $B-(N-1)$——这正好是独立回路的个数，
这就是第 2 条。KVL 侧的同一条事实写成 $\vect{A}\vect{A}_a^{\mathsf{T}}=\vect{0}$（回路矩阵乘关联矩阵转置为零）。}

\subhdr{1.3 逐页核对 p7--p10：$M$ 有公式，$L$ 没有}

四页计数练习的答案都是印好的，本讲义做的是\textbf{复核}，并把「哪些数能算、哪些数只能数」分开：

\par\vspace{2pt}
\noindent\begin{tabular}{@{}p{0.9cm}p{0.8cm}p{0.9cm}p{1.3cm}p{2.0cm}p{1.1cm}p{7.0cm}@{}}
\tabtop
\thead{页} & \thead{$N$} & \thead{$B$} & \thead{$M$（课件）} & \thead{$B-N+1$} & \thead{$L$（课件）} & \thead{说明}
\\
\tabmid
p7 & 2 & 3 & 2 & $3-2+1=2$ & 3 & 三条支路并接在同一对节点之间；任取两条就围成一个回路，共 $C_3^2=3$ 个
\\
p8 & 2 & 4 & 3 & $4-2+1=3$ & 6 & 四条支路并接在同一对节点之间；任取两条围成一个回路，共 $C_4^2=6$ 个
\\
p9 & 4 & 6 & 3 & $6-4+1=3$ & 7 & 桥式结构，$M$ 与公式吻合；$L=7$ 的拓扑依据见下方说明
\\
p10 & 3 & 5 & 3 & $5-3+1=3$ & 6 & 三个网孔串成一串（左、中、右），只有相邻两个才能合并成大回路，$3+2+1=6$
\tabbot
\end{tabular}

\textbf{结论有两层，这比四个答案本身有用：}
（1）$M$ 只依赖 $N$ 与 $B$，有公式，可以算；
（2）$L$ 依赖「网孔怎么连」，没有只含 $N,B$ 的公式——p8、p10 的 $M$ 都是 3 且 $L$ 都是 6，
p9 的 $M$ 也是 3 而 $L=7$，三个电路同 $M$ 不同 $L$，这就证明 $L$ 不是 $M$ 的函数。
所以\textbf{列方程时永远数网孔 $M$，不数回路 $L$}。

\aside{关于 p9 的 $L=7$：扫描图上能确认的是 $N=4$、$B=6$、$M=3$，
而 $6$ 条支路、$4$ 个节点要凑出 $7$ 条简单回路，最典型的情形是「四点全连」（$K_4$：
$4$ 个三角形 $+$ $3$ 个四边形）——课件那张图是不是 $K_4$，文字层里看不出来。
原图在文字层上只留下三个 $+/-$ 记号与四个电阻框的位置，桥臂的接法无法唯一复原，
故 $L=7$ 这一项本讲义只备案、不复核，标 [待确认]。}

\subhdr{1.4 用 p49 验一遍：为什么是「5 equations」}

p49 的例题电路有 3 个节点、5 条支路，页面上直接印了「How many equations are needed? 5 equations」。
按 §1.2：$N-1=2$ 条 KCL，$M=B-N+1=5-3+1=3$ 条 KVL，合计 $2+3=5=B$，正好等于未知数个数。
这一步在 §8.1 会落到每一行方程上；p52 那句「Only two of them are independent」
（三个节点方程里只有两条独立）与「We have two equations. We need three more.」
（还差三条 KVL）就是同一条定理的两个半句——课件把它拆在例题里讲，没有在 p48 说过。

\subhdr{1.5 三张图逐个认一遍：数之前先圈「三条以上支路相交的点」}

光有公式不行，得会认图。p7、p8、p10 这三页正好是三种典型结构：

\par\vspace{2pt}
\noindent\begin{tabular}{@{}p{1.1cm}p{6.3cm}p{3.0cm}p{4.4cm}@{}}
\tabtop
\thead{页} & \thead{图上有什么（按扫描图读）} & \thead{数出来的} & \thead{为什么}
\\
\tabmid
p7 & 三条支路并接在同一对节点之间：左支路一只电阻串 $5\,\mathrm{V}$ 电压源（$+$ 在上），
中支路一只 $1\,\mathrm{A}$ 电流源（箭头向上），右支路一只电阻 $R$ & $N=2$，$B=3$，
$M=2$，$L=3$ & 三条支路的上端接同一条导线、下端接另一条导线，两个交点都汇了三条支路，正好是节点；
任取两条支路围成一个回路，$C_3^2=3$
\\
p8 & 四条支路并接在同一对节点之间（图上只剩三个 $+/-$ 记号，元件细节在文字层里没有留下） &
$N=2$，$B=4$，$M=3$，$L=6$ & 同上；$C_4^2=6$
\\
p10 & 五条支路、三个节点：左右各一条竖支路，上边一条含 $R_5$ 的横支路，中间两条竖支路 $R_3$、$R_4$；
两处 $+/-$ 记号说明其中有电压源 & $N=3$，$B=5$，$M=3$，$L=6$ & 三个网孔串成一串
\tabbot
\end{tabular}

\textbf{三页的共同教训有两层}：
（1）数的方法永远一样——\textbf{先圈「三条以上支路相交的点」（那就是节点），再数两个节点之间的一段（那就是支路）}；
（2）p7 的中支路是\textbf{电流源}，它照样算一条支路——\textbf{拓扑数与元件类型无关}。
这也是 §1.2 的定理只含 $N$ 与 $B$ 的原因：KCL 与 KVL 都不看元件是什么（§2.2、§3.4）。

\warnbox{\textbf{一个容易多数的地方}：两个元件\textbf{首尾相接}的那一点不算节点（p6 的原文是
「three or more elements」）。p7 的左支路里「电阻与 $5\,\mathrm{V}$ 源相接」的那一点只有两个元件，
不是节点；若把它当节点，$N$ 就数成 3，方程数也跟着多一条。判断办法：
\textbf{这个点上有几条支路汇合？}两条就不是节点，三条及以上才是。}

\subhdr{1.6 一句补白：$M=B-N+1$ 对任何连通电路都对，「网孔」才要求平面}

这条定理有个容易误解的地方，要点明：$M=B-N+1$ 里的 $M$ 严格说是\textbf{独立回路的个数}
（回路空间的维数），它对\textbf{任何连通电路}都成立，不管能不能画在平面上；
而「网孔」这个词只对\textbf{平面电路}有意义——必须能画成「导线只在节点处相交、不跨接」的样子。
所以：
\begin{itemize}
\item 平面电路：网孔就是独立回路，数网孔最方便（p7--p10、p49 都是这种）；
\item 非平面电路（例如把立方体的棱当支路的电路）：没有「网孔」这个对象，
但独立回路仍然是 $M=B-N+1$ 个，此时要靠「先取一棵树，每加一条连支成一个回路」的办法来找（本册不展开）。
\end{itemize}

\aside{本册所有电路都能画平（p7--p10、p37--p42、p49--p56 全是平面图），
所以 §7.3 的实用数法「每个网孔取一条」用得上。
但公式本身不依赖平面性——这决定了它后面能直接搬到回路电流法里去用。}

%=============================================================
\hdr{§2}{KCL：从「电荷不堆积」到「代数和为零」}

\subhdr{2.1 课件 p11--p13 的三页，其实是「一份提纲 + 一句陈述」}

课件 p11 用五项做提纲，五项下面\textbf{都是空的}。原文这五项依次是
「For:」「Content:」「Examples:」「Other expression:」「Generalized node:」，
前面各带一个圆圈编号（本讲义一律改成全角括号编号，见 §0.3）；
p12 只把第一项「For:」标红并画了一个节点电流图（一个节点 A、五条支路、五支箭头），
p13 给出陈述：

\callout{原文（p13）：\quad「KCL states that the algebraic sum of the current entering a node
( or a closed boundary) is zero.」\quad 并附 $\Sigma i=0$、$\Sigma I=0$ 两式。\\
译：\textbf{流入一个节点（或一个闭合边界）的电流的\emph{代数和}为零}。}

这句话本身没错，但它有两个洞，两个洞都不小：
（1）「algebraic（代数的）」二字要求每个电流带正负号，\textbf{可课件没说符号怎么定}——
直到 p14 才用一个例子（$I_1+I_3-I_2-I_4-I_5=0$）把答案暗示出来；
（2）\textbf{为什么会有这条定律}，一页也没有。下面两小节各补一个洞。

\subhdr{2.2 为什么成立：节点没有地方存电荷}

把要用的前提说全：\textbf{集中参数假设}——所有元件都集中在「点」上，连接线是理想导线，
节点的几何尺寸远小于电路工作频率对应的波长，因此\textbf{节点的电容可以忽略，电荷不能在节点上堆积}。

在这个前提下，画一个很小的闭合面把节点包住，对这个面用电荷守恒：
单位时间进入这个面的电荷量，必须等于单位时间离开的电荷量；若不等，
多出来的电荷只能堆在节点上，节点电位就会不断升高，这与「电位是有限值」矛盾。
写成式子就是对每条穿过该面的支路把电流带符号相加为零——这正是 KCL。

\aside{KCL 的成立\textbf{不依赖元件的伏安关系}：电阻、电感、电容、电源都一样，
因为它来自「电荷守恒 + 集中参数」，与元件是线性还是非线性无关。
课件 p4 那句「Kirchhoff's laws can be used for the analysis of DC circuits,
AC circuits and nonlinear circuits」正是这条性质，只是没给理由。}

\subhdr{2.3 符号规则：三种等价说法，一题之内只能用一种}

\warnbox{\textbf{（甲）流入取正、流出取负，代数和为零}：$\sum I_{\text{入}}-\sum I_{\text{出}}=0$；
\textbf{（乙）流入之和 = 流出之和}：$\sum I_{\text{入}}=\sum I_{\text{出}}$；
\textbf{（丙）流出取正、流入取负}：与（甲）只差一个负号，同样对。\\
三种写法\textbf{互相等价}（甲与乙只差移项，甲与丙只差乘 $-1$），选一种用到底即可。
\textbf{禁忌}：同一题里这条按甲、那条按乙，正负号必然乱。}

p14 的方程逐项对照（图上的箭头方向是按 p12 画的，本讲义按 §0.3 的约定读）：

\par\vspace{2pt}
\noindent\begin{tabular}{@{}p{2.0cm}p{4.2cm}p{3.0cm}p{6.0cm}@{}}
\tabtop
\thead{电流} & \thead{箭头相对节点 A} & \thead{按（甲）取的符号} & \thead{说明}
\\
\tabmid
$I_1$ & 指向 A（流入） & $+I_1$ & 左上支路，箭头指进节点
\\
$I_3$ & 指向 A（流入） & $+I_3$ & 右侧支路，箭头指进节点
\\
$I_2$ & 背离 A（流出） & $-I_2$ & 右上支路，箭头离开节点
\\
$I_4$ & 背离 A（流出） & $-I_4$ & 右下支路
\\
$I_5$ & 背离 A（流出） & $-I_5$ & 左下支路
\tabbot
\end{tabular}

代入得 $I_1+I_3-I_2-I_4-I_5=0$，与 p14 的第一式完全一致；
移项即 $-I_2-I_4-I_5=-I_1-I_3$，两边乘 $-1$ 得 $I_2+I_4+I_5=I_1+I_3$，
这就是 p14 那句「or $I_1+I_3=I_2+I_4+I_5$」——\textbf{课件的两个写法是同一个方程的两种形式（甲）与（乙）}，
不是两条定律。看懂这一点，「符号怎么定」的问题就消掉了一半：先定一种，然后\textbf{整式一起乘 $-1$ 也不改变结论}。

\subhdr{2.4 广义节点（p15--p17）：为什么可以「把一整块电路当一个节点」}

课件 p13 的括号里留了一句「( or a closed boundary)」，p15--p17 用三页演示它的意思：
p15 是三条支路穿过包围一块电路的区域，p16 是三角形电阻网络，p17 是四边形电阻网络，
每页下方都印着 \texttt{closed boundary} 与第五项「Generalized node:」的红色高亮。

\warnbox{\textbf{广义节点（KCL 的推广形式）}\quad
任取一个把若干节点和支路包在里面的\textbf{闭合面}（closed boundary），
则\textbf{所有穿过该闭合面的支路}电流的代数和为零；面内支路的电流不出现在方程里。}

\textbf{为什么成立（证明只有一行）}：把面内每一个节点的 KCL 方程相加。
面内每一条支路连接两个面内节点，它在这两个节点的方程里各出现一次、符号相反，
所以\textbf{全部两两抵消}；剩下的恰好是穿出面的那些支路电流，代数和为零。
这也顺便说明了第二条注意事项：\textbf{只能统计穿过边界的支路}。

\aside{p16 的三条外引电流（$I_1,I_2,I_3$）与 p17 的四条（$I_1\sim I_4$）就是「穿过闭合面的支路」。
p15 与 p16/p17 的图都只画了方向、没写方程，练习时按 §2.3 选一种符号规则，
把「流入面」取正、「流出面」取负，三项或四项相加为零即可。}

\subhdr{2.5 一句话记住 KCL}

\textbf{KCL 只问「接了几条线、每条线是往里还是往外」，不问线上是什么元件。}
所以它的用法永远是三步：认定节点（或闭合面）$\to$ 给每条支路定参考方向 $\to$
按 §2.3 的规则带符号相加。

\subhdr{2.6 三种最常见的误用，以及各自的自检办法}

\par\vspace{2pt}
\noindent\begin{tabular}{@{}p{0.8cm}p{6.0cm}p{7.9cm}@{}}
\tabtop
\thead{序} & \thead{误用} & \thead{自检办法}
\\
\tabmid
（1） & 把「两个元件首尾相接」的接点当节点，$N$ 数多了 &
数这个点上有几条支路汇合；两条就不是节点（§1.5）
\\
（2） & 同一题里一会儿按「流入为正」、一会儿按「流出为正」 &
在图上给每条支路画一支箭头，然后逐项念「指向节点 $\to$ 正、背离节点 $\to$ 负」，
念完一遍检查有没有中途换约定
\\
（3） & 用广义节点时把面内支路也写进方程 &
只统计\textbf{穿过边界}的支路；面内支路在求和时已经一正一负抵消掉了（§2.4 的证明）
\tabbot
\end{tabular}

\warnbox{\textbf{一句可以背下来的话}：KCL 的方程里，\textbf{每一项的符号只由「那支箭头的方向」
与「你把谁定成正」这两件事决定}，与电流的大小、元件的种类、电压的高低都没有关系。
所以 KCL 是全册\textbf{最好写}的方程——写错一定是「方向记混」，不可能是「物理不懂」。}

%=============================================================
\hdr{§3}{KVL：从「绕一圈回到原电位」到「代数和为零」}

\subhdr{3.1 课件 p18--p20：先定对象，再定方向，最后才给陈述}

三页的顺序是有讲究的，不能跳：
p18 用一行 \texttt{Loop\ \ Mesh} 先\textbf{定对象}（有四个电阻首尾相接围成的那个闭合路径）；
p19 在同一个网孔上加一条回环箭头并标 \texttt{Loop direction}，\textbf{定绕行方向}；
p20 才给出陈述：

\callout{原文（p20）：\quad「KVL states that the algebraic sum of all voltages around a
closed path (or loop) is zero.」\quad 并附 $\Sigma u=0$、$\Sigma U=0$。\\
译：\textbf{沿任一闭合路径（回路）绕行一周，各段电压的代数和为零}。}

这里「algebraic」二字的含义与 KCL 一样含糊：\textbf{哪一段取正、哪一段取负}，p20 一个字没写。
p21--p23 用 E1 演示了一遍（先标四个电阻的极性，再标绕行方向，最后写出 $U_1-U_2+U_3-U_4=0$），
但\textbf{中间那一步判断规则始终没写出来}。§3.2 把它写出来，§3.3 与 §4 各验证一次。

\subhdr{3.2 统一符号规则（本讲义的核心补法）}

\warnbox{\textbf{规则（本讲义用这条，全册一致）}\quad 给回路上每个元件配一个\emph{箭头}：
\textbf{电压源的箭头从 $+$ 指向 $-$}；\textbf{电阻的箭头取电流的参考方向}。
然后沿回路绕行一周：\textbf{元件的箭头与绕行方向一致，该项取正；相反，取负}。代数和为零。\\
\textbf{等价的端子说法}：沿绕行方向从 $+$ 到 $-$ 取正，从 $-$ 到 $+$ 取负。}

\aside{两种说法为什么等价：电压源的箭头是「$+\to-$」（电位降的方向），
所以「箭头与绕行同向」就是「从 $+$ 走到 $-$」。
电阻没有固定极性，它的「$+$ 端」只能靠电流定义——电流流入的那端就是电位高的一端，
于是「箭头与绕行同向」就是「顺着电流走」，这一段的压降取 $+RI$。}
 
\subhdr{3.3 拿 E1 走一遍（p21--p23）：四段四个符号，正好对上课件的结果}

p21 给四个电阻标了极性，p23 在中央加了一个\textbf{顺时针}的回环箭头（Loop direction），
并在四个元件外侧画了四个红色箭头。先把 p21/p23 上的事实列成表：

\par\vspace{2pt}
\noindent\begin{tabular}{@{}p{1.5cm}p{1.3cm}p{2.8cm}p{2.9cm}p{2.4cm}p{4.2cm}@{}}
\tabtop
\thead{段} & \thead{电压} & \thead{极性（$+$ 端在哪）} & \thead{元件箭头指向} & \thead{绕行方向} & \thead{符号}
\\
\tabmid
上边 & $U_1$ & 左端 $+$、右端 $-$ & 向右（$+\to-$） & 向右 & 一致，取 $+U_1$
\\
右边 & $U_2$ & 下端 $+$、上端 $-$ & 向上（$+\to-$） & 向下 & 相反，取 $-U_2$
\\
下边 & $U_3$ & 右端 $+$、左端 $-$ & 向左（$+\to-$） & 向左 & 一致，取 $+U_3$
\\
左边 & $U_4$ & 上端 $+$、下端 $-$ & 向下（$+\to-$） & 向上 & 相反，取 $-U_4$
\tabbot
\end{tabular}

四段相加：
\[ U_1-U_2+U_3-U_4=0 , \]
与 p23 下方印的那一行\textbf{逐项一致}。这一步值得停一下，因为初学者最容易问的是
「凭什么上、下取正、左、右取负」——答案就在表里：\textbf{四个元件箭头本身绕成一个逆时针的圈，
而绕行方向是顺时针}，所以每遇到一个「与绕行相反」的元件（右边的 $U_2$、左边的 $U_4$）就取负。

\aside{四个元件箭头的极性也可以直接从 p21 的四个 $+/-$ 记号读出来：
$U_1$ 的 $+$ 在左上、$U_2$ 的 $+$ 在右下、$U_3$ 的 $+$ 在右下、$U_4$ 的 $+$ 在左上。
本表的「极性」一列就是从这四个记号抄下来的，没有改。}

\subhdr{3.4 为什么这条规则一定对：电位是单值的}

规则不是背出来的，它是\textbf{「电位单值」这一条}的记账方式：
从回路某点出发，沿绕行方向走过第 $k$ 段时电位的变化量记为 $\Delta V_k$
（电位降低记为正的「降」），绕一圈回到出发点，总变化必须为零：
\[ \sum_k \Delta V_k = 0 . \]
而每一段的「降」只有两种算法：顺电流方向过电阻，降 $=+R_kI_k$（逆着电流则为 $-R_kI_k$）；
从 $+$ 极走到 $-$ 极过电源，降 $=+U_S$（从 $-$ 走到 $+$ 则为 $-U_S$）。
把这两种算法代进上式，就得到 §3.2 的规则——\textbf{它只是「电位单值」的记账语言}。

\aside{注意 KVL 与 KCL 的对称性：KCL 来自「电荷守恒 + 集中参数」，KVL 来自「电位单值」。
两条都不依赖元件的伏安关系，所以课件 p4 说它们对非线性电路也成立。
真正用到元件性质的是最后一步「$U_R=RI$」——这一条把电压与电流联系起来，
KCL+KVL 才能解出数。}

\subhdr{3.5 三个推论：回路的选法为什么可以随便、假想支路为什么合法}

\textbf{推论一（两点电压与路径无关）}：同一电路中，从 $P$ 到 $Q$ 沿两条不同路径走，
电压降必须相等——因为两条路径合起来就是一个闭合回路，KVL 说这条回路的电压和为零，
于是「沿路径 1 的降」$=$ 「沿路径 2 的降」。这条推论是 §5 的全部基础：
\textbf{只要两个端点的电压起了名字，这一段就能当元件}。

\textbf{推论二（无元件的导线不影响方程）}：一段只有导线、没有元件的路径，它的电压降为零，
写进 KVL 就是 $+0$。所以回路可以沿着导线任意「拐弯」，只要首尾相接。

\textbf{推论三（开路支路可以忽略，但不能当它不存在）}：p21--p23 的 E1 在四个角上各引出一段
\emph{开路}短支路（p21 左上角还标了粉框 E1）。开路支路的电流为零，写进 KVL 没有任何贡献，
所以四个电阻的方程照样成立；但这段支路的存在说明\textbf{这块网络还连着别处}，
后面若把它接进更大的电路，这四个电压 $U_1\sim U_4$ 就会跟着变——这一点在 1.10 等效电源定理里会用上。

\aside{推论一反过来读也成立：若沿两条路径算出的电压降不相等，那一定是算错了，
不必去怀疑 KVL。这是做 KVL 题最快的一种自检——同一对端子换一条路径再走一遍，
两个结果必须相同（§6.4 里 $U_{ab}$ 就用了这个办法：既从右上角算、也从下方的 10 V 源算）。}

%=============================================================
\hdr{§4}{E2（p24--p27）：电压源与电阻的「两个箭头」不是一回事}

\subhdr{4.1 电路与方向}

p24 摆出电路：左边一条支路（$U_{S1}$ 在上、$R_1$ 在下），右边一条支路（$U_{S2}$ 在上、$R_2$ 在下），
两条支路经上下两条水平导线接成单回路；两台电压源都是\textbf{上端 $+$、下端 $-$}。
p25 加上一圈\textbf{顺时针}的弧箭头（Loop direction）；
p26 再补四个红色箭头：$U_{S1}$ 旁向下、$R_1$ 旁向上、$U_{S2}$ 旁向下、$R_2$ 旁向下。
p27 给出结果：

\[ U_{S2}+R_2I_2+R_1I_1-U_{S1}=0 . \]

\subhdr{4.2 用 §3.2 的规则逐项复算}

绕行方向取顺时针，即「沿上边向右 $\to$ 沿右支路向下 $\to$ 沿下边向左 $\to$ 沿左支路向上」：

\par\vspace{2pt}
\noindent\begin{tabular}{@{}p{1.9cm}p{2.0cm}p{3.4cm}p{2.5cm}p{2.3cm}p{3.0cm}@{}}
\tabtop
\thead{段} & \thead{元件} & \thead{箭头指向} & \thead{绕行方向} & \thead{符号} & \thead{理由}
\\
\tabmid
右支路 & $U_{S2}$ & 向下（$+\to-$） & 向下 & $+U_{S2}$ & 一致
\\
右支路 & $R_2$ & 向下（$I_2$） & 向下 & $+R_2I_2$ & 顺着电流
\\
左支路 & $R_1$ & 向上（$I_1$） & 向上 & $+R_1I_1$ & 顺着电流
\\
左支路 & $U_{S1}$ & 向下（$+\to-$） & 向上 & $-U_{S1}$ & 相反
\tabbot
\end{tabular}

四项相加即 $U_{S2}+R_2I_2+R_1I_1-U_{S1}=0$，与 p27 \textbf{完全相同}。
\textbf{这是同一套规则第二次独立通过验证}（第一次是 §3.3 的 E1）——
两页的电路、方向、元件都不同，规则却是同一条，所以它才配叫「规则」。

\subhdr{4.3 两个必须点名的细节}

\gapitem{（1）四个红色箭头分两类}{左支路里「$U_{S1}$ 旁向下」和「$R_1$ 旁向上」方向相反，
不是画错了：前者是\textbf{电压参考方向}（$+\to-$，由极性决定），后者是\textbf{电流参考方向}（由解题者自己定）。
同一条支路上这两类箭头可以同向、也可以反向，决定于「电流实际从电源的哪一端流出」。
判断时\textbf{不要看它画在哪个元件旁就当同类}。}

\gapitem{（2）p27 的末项印成 $U_{1S}$}{原文照录：\texttt{US2+ R2I2 + R1I1－U1S＝0}，
末项的下标顺序是反的。同一符号在 p24--p26 与 §3.3 的 E1 里都写成 $U_{S1}$，
本讲义按 $U_{S1}$ 使用；这是印刷笔误，不影响方程结构。}

\subhdr{4.4 一条自检：把绕行方向反过来，整式应变号}

如果改取逆时针绕行，四段的「一致/相反」全部翻转，得
$-U_{S2}-R_2I_2-R_1I_1+U_{S1}=0$，与原式只差一个 $-1$——\textbf{两种绕行方向给出同一个方程}。
这不是巧合：绕行方向反了，$-\sum \Delta V_k=0$ 与 $\sum \Delta V_k=0$ 同解。
所以\textbf{绕行方向怎么取都可以，取好以后全回路必须一致}；
真正不能乱的是「同一条支路上元件的箭头」与「绕行方向」谁跟谁同向——那是 §3.2 管的。

\aside{移项后得到 \ $U_{S1}=U_{S2}+R_2I_2+R_1I_1$。
这个写法在 §8 的支路电流法里会反复出现：\textbf{每一条 KVL 都能读成「这个网孔里电源的电动势
$=$ 各段电阻压降之和」}（沿同一个绕行方向）。p53 的三条方程全是这个形状，
只是把电阻项按支路分别写出来了。}

%=============================================================
\hdr{§5}{广义回路与端口电压（p28--p32）：$U=U_S+RI$ 是怎么来的}

\subhdr{5.1 p28--p29：把「一个电源 + 一个电阻」接到任意网络上}

p28 的图是本节的主角：一条竖直支路上，电压源 $U_S$（上端 $+$、下端 $-$）与电阻 $R$ 串联，
这条支路向右接到一个\textbf{没有任何标注的方块}（外部电路）；向左，从支路的上下两端各引出一条短导线，
两线之间标一个电压 $U$（上端 $+$、下端 $-$）。
p29 用一圈蓝色粗线把「$U$ 所在的那条\emph{假想}竖直线 $+$ 上导线 $+$ $U_S/R$ 支路 $+$ 下导线」圈成一个
闭合回路，框外标 \texttt{Imaginary loop}。

这一步是本册最有用的一个观念：\textbf{KVL 只要求「一条闭合路径」，不要求路径上都是真元件}。
一条两端电压已经起好名字（这里是 $U$）的「假想支路」，完全可以充当回路的一段。
课件把这件事列为第五项「Generalized loop:」，与 KCL 的「Generalized node」并列。

\aside{「假想支路」的电压 $U$ 不是算出来的，是\textbf{给定}的——它就是这对端子的电压。
后面 1.10 等效电源定理（[C1-4]）要问的是「这个方块能不能换成一个电源串电阻还不改变 $U$」，
起点正是这一页：先把「从端口看进去」这件事说清楚。本册不展开，只记一句：
\textbf{端口电压与方块内部结构无关，只要还是同一对端子}。}

\subhdr{5.2 p30--p31：方向先定，再列式}

p30 补上两个方向，p31 又把电压的参考方向补齐（页下方印着 「Find the direction of all the voltage」）：
\begin{enumerate}
\item \textbf{绕行方向}：回环中央一支青色弧箭头，按「上边向右、右侧向下、下边向左、左侧向上」绕——顺时针；
\item \textbf{电流参考方向}：支路里一支蓝色箭头由 $U_S$ 下端穿 $R$ 向下，标 $I$，即\textbf{$I$ 向下}；
\item \textbf{电压参考方向}：$U$ 旁与 $U_S$ 旁各一支红色向下箭头（$+\to-$）。
\end{enumerate}

注意第 3 条：$U$ 的 $+$ 在上端、$U_S$ 的 $+$ 在上端，两支红箭头都向下，所以\textbf{这两个电压的参考方向是同向的}。
这一点决定了下式的符号，也是 §6 里「反一个电流就换一个符号」的由来。

\subhdr{5.3 p32：$U_S+RI-U=0$ 的逐项复算}

沿用 §3.2 的规则，绕行方向顺时针（左段向上、右段向下）：

\par\vspace{2pt}
\noindent\begin{tabular}{@{}p{2.1cm}p{1.4cm}p{3.6cm}p{2.6cm}p{2.0cm}p{2.6cm}@{}}
\tabtop
\thead{段} & \thead{量} & \thead{箭头指向} & \thead{绕行方向} & \thead{符号} & \thead{说明}
\\
\tabmid
左（假想支路） & $U$ & 向下（$+\to-$） & 向上 & $-U$ & 相反
\\
右（$U_S$） & $U_S$ & 向下（$+\to-$） & 向下 & $+U_S$ & 一致
\\
右（$R$） & $R$ & 向下（$I$ 向下） & 向下 & $+RI$ & 顺着电流
\tabbot
\end{tabular}

三项相加为零：
\[ U_S+RI-U=0 \quad\Longrightarrow\quad U=U_S+RI , \]
与 p32 印的两行\textbf{完全一致}（第一行是方程，第二行是移项后的结果）。
\textbf{这是同一套规则第三次通过验证}。

\aside{移项后那行 $U=U_S+RI$ 的读法值得记：\textbf{端口电压 = 电源电压 + 电阻上的压降}。
在「电源与负载同向串在这条支路上」的常见画法里，$U$ 总是比 $U_S$ 大——因为电流在 $R$ 上又掉了
一截 $RI$。§6.1 把电流反过来以后，这一截就变成「抵掉」，于是 $U=U_S-RI$。}

\subhdr{5.4 p37：已知电流求端口电压（$U_{AB}=12\,\mathrm{V}$）}

p37 是同一个观念的数值版。图上只有一句话 \texttt{Find UAB}，再没有别的要求。
能确认的事实是：$A$、$B$ 是两个端子（$A$ 在上、$B$ 在下），$U_{AB}$ 的参考方向是一支由 $A$ 指向 $B$ 的箭头
（所以 $U_{AB}=V_A-V_B$）；$A$、$B$ 之间那条支路上，上面是 $2\,\Omega$、下面是 $40\,\mathrm{V}$ 电压源（上端 $+$、下端 $-$），
支路旁标 $14\,\mathrm{A}$ 与一支\textbf{向上}的箭头（电流参考方向向上）。

既然\textbf{电流已知}，就不必去解右边的 $2\,\Omega/20\,\mathrm{V}$ 与 $4\,\Omega$ 两条支路——
沿已知的那一条走一遍即可（走的方向：$A\to B$，即\textbf{向下}）：

\par\vspace{2pt}
\noindent\begin{tabular}{@{}p{2.5cm}p{3.1cm}p{3.2cm}p{2.2cm}p{4.1cm}@{}}
\tabtop
\thead{元件} & \thead{箭头指向} & \thead{沿 $A\to B$ 走} & \thead{电位变化} & \thead{依据}
\\
\tabmid
$2\,\Omega$ & 向上（$I=14\,\mathrm{A}$） & 逆着电流走 & 升 $2\times14=28\,\mathrm{V}$ & 逆电流方向电位升高 $RI$
\\
$40\,\mathrm{V}$ 源 & 向下（$+\to-$） & 顺着 $+\to-$ 走 & 降 $40\,\mathrm{V}$ & 从 $+$ 走到 $-$，电位降低 $U_S$
\tabbot
\end{tabular}

从 $A$ 出发、到 $B$ 结束，电位的净变化是「先升 28、再降 40」：
\[ V_B-V_A=28-40=-12\,\mathrm{V} \quad\Longrightarrow\quad U_{AB}=V_A-V_B=+12\,\mathrm{V} . \]

\aside{换个写法更快：$U_{AB}$ 就是「沿 $A\to B$ 走时各段的电位降之和」，即
$U_{AB}=\sum(\text{降})-\sum(\text{升})=40-28=12\,\mathrm{V}$。
与 §5.3 的 $U_S+RI-U=0$ 是同一件事的两种记法——那边把未知量留在左边，这边把未知量直接解出来。}

\gapitem{（本页的一处不自洽，标 [待确认]）}{若把右边两条支路也当作与左边那条\textbf{并联}在同一对母线上，
则按 $40\,\mathrm{V}/2\,\Omega$、$20\,\mathrm{V}/2\,\Omega$、$4\,\Omega$ 三条支路解整个网络，
左边那条支路的电流应算得 $8\,\mathrm{A}$，而不是图上标的 $14\,\mathrm{A}$。
这说明 $14\,\mathrm{A}$ 是\textbf{作为已知条件外加上去的}，不是从这个网络解出来的。
本讲义按「14 A 已知」算 $U_{AB}=12\,\mathrm{V}$；右边两条支路究竟怎么接，原图分辨率下不能唯一确定，
标 [待确认]，不作为考点。}

\subhdr{5.5 端口电压的三种约定：换教材之前先看一眼 $+/-$ 记号}

同一个端口电压，不同教材的画法不同，算出来的数会差一个负号。本册出现的三种：

\par\vspace{2pt}
\noindent\begin{tabular}{@{}p{1.0cm}p{5.6cm}p{4.0cm}p{4.4cm}@{}}
\tabtop
\thead{约定} & \thead{画法} & \thead{数学含义} & \thead{本册在哪出现}
\\
\tabmid
A & 上端标 $+$、下端标 $-$，参考箭头向下（由 $+$ 指向 $-$） & $U=V_{\text{上}}-V_{\text{下}}$ &
p28--p32 的 $U$、p37 的 $U_{AB}$、p42 的 $U_{ab}$
\\
B & 上端标 $-$、下端标 $+$ & $U=V_{\text{下}}-V_{\text{上}}$（即约定 A 那个 $U$ 的相反数） &
p36 的 $U_3$、$U_4$ 两个子图（正极在下）
\\
C & 写成 $U_{ab}$ 且 $a$ 在上 & 与 A 相同（$U_{ab}=V_a-V_b$） & p38--p42
\tabbot
\end{tabular}

判据只有一条：\textbf{看题干给的箭头与 $+/-$ 记号，两者必须自洽}（箭头应当从 $+$ 指向 $-$）。
若图上有箭头没记号，以箭头为准；有记号没箭头，以记号为准；
\textbf{两者都没有时}，按「从上到下」取约定 A，并在答案里写明「取上端为正」。

%=============================================================
\hdr{§6}{电流参考方向反过来（p33--p36）：一个负号引起的连锁}

\subhdr{6.1 p33--p35：$U=U_S-RI$ 的来历}

p33 只在页顶加了一句浅绿的问句：「What if the current is in the opposite direction?」
其余什么都没改——\textbf{电压源的极性、$U$ 的极性、绕行方向全都保持原样}，只把 $I$ 的蓝色箭头由向下改成\textbf{向上}
（p34 补齐外部方块的回路电流、p35 给出结果）：

\[ U_S-RI-U=0 \quad\Longrightarrow\quad U=U_S-RI . \]

逐项看，四个量里只有 $R$ 那一项换了符号，因为只有它的箭头（电流参考方向）反了：

\par\vspace{2pt}
\noindent\begin{tabular}{@{}p{2.3cm}p{3.5cm}p{3.3cm}p{2.4cm}p{2.8cm}@{}}
\tabtop
\thead{量} & \thead{p32（$I$ 向下）} & \thead{p35（$I$ 向上）} & \thead{箭头与绕行} & \thead{符号}
\\
\tabmid
$U$ & 箭头向下 & 箭头向下（未变） & 相反 & $-U$（两页相同）
\\
$U_S$ & 箭头向下 & 箭头向下（未变） & 一致 & $+U_S$（两页相同）
\\
$R$ & 箭头向下 & 箭头\textbf{向上} & 由一致变相反 & $+RI\Rightarrow-RI$
\tabbot
\end{tabular}

\subhdr{6.2 判据：两项各看各的，不要混}

\warnbox{\textbf{两句话判完（本册全部符号问题都能归结到这两条）}\\
（甲）\textbf{电阻项}：沿绕行方向走时，\emph{顺着}电流参考方向取 $+R_kI_k$，\emph{逆着}取 $-R_kI_k$。\\
（乙）\textbf{电源项}：沿绕行方向走时，从 $+$ 走到 $-$ 取 $+U_S$，从 $-$ 走到 $+$ 取 $-U_S$。\\
两条互不影响：电流反了只动（甲），电源画反了只动（乙），两者都反则两项同时变号。}

这一条判据直接解释了为什么 p32 与 p35 只差一个符号，也解释了 p23 的四个电阻为什么两个正两个负（§3.3）。
\aside{记住「$+R_kI_k$ 的 $+$ 不是电流方向，是\emph{走的方向与电流方向一致}」。
初学者常把「电流向上所以取正」当成规则，于是同一个元件换个绕行方向就出错。
正确的写法永远是「两个方向比较」，绕行方向是分母。}

\subhdr{6.3 p36 四小题：$U_1\sim U_4$ 与四个式子的对应表}

p36 一页塞了四个子图，页顶一句 \texttt{Find U1，U2，U3，U4}，页下四行待填式子：
「(1)$U\ \square=U_S+IR$」「(2)$U\ \square=-U_S-IR$」「(3)$U\ \square=U_S-IR$」「(4)$U\ \square=-U_S+IR$」，
方框是留给 $U_1\sim U_4$ 的位置。\textbf{课件没有给答案}，下面补上。

四个子图的拓扑完全一样：一根竖直支路，\textbf{上端是电阻 $R$、下端是电压源 $U_S$}，
支路上、下两端各引一段短导线作端子；$U_k$ 量的就是这对端子之间的电压（上端标 $+$、下端标 $-$），
电流箭头画在 $R$ 上方。四个子图只差两件事：\textbf{电压源的正极在上还是在下}、\textbf{电流箭头向上还是向下}。
把这两件事按（甲）（乙）两条判据组合，只有四种情况：

\par\vspace{2pt}
\noindent\begin{tabular}{@{}p{1.2cm}p{2.4cm}p{1.8cm}p{3.6cm}p{3.2cm}p{2.8cm}@{}}
\tabtop
\thead{子图} & \thead{$U_S$ 正极} & \thead{$I$ 方向} & \thead{$R$ 项（从 $+$ 端子走到 $-$ 端子）} & \thead{$U_S$ 项} & \thead{结果}
\\
\tabmid
$U_1$ & 在上 & 向上 & 逆电流，电位升 $RI$，取 $+RI$ & 从 $+$ 到 $-$，取 $-U_S$ & $U_S-RI$
\\
$U_2$ & 在上 & 向下 & 顺电流，电位降 $RI$，取 $-RI$ & 从 $+$ 到 $-$，取 $-U_S$ & $U_S+RI$
\\
$U_3$ & 在下 & 向上 & 逆电流，取 $+RI$ & 从 $-$ 到 $+$，取 $+U_S$ & $-U_S-RI$
\\
$U_4$ & 在下 & 向下 & 顺电流，取 $-RI$ & 从 $-$ 到 $+$，取 $+U_S$ & $-U_S+RI$
\tabbot
\end{tabular}

于是对应关系是\textbf{一个错位}，不是逐行照抄：
\[ U_1\ \longleftrightarrow\ (3)U=U_S-IR,\qquad
   U_2\ \longleftrightarrow\ (1)U=U_S+IR, \]
\[ U_3\ \longleftrightarrow\ (2)U=-U_S-IR,\qquad
   U_4\ \longleftrightarrow\ (4)U=-U_S+IR . \]

\aside{代一组数自检：取 $U_S=10\,\mathrm{V}$、$R=2\,\Omega$、$I=1\,\mathrm{A}$，
则 $U_1=8$、$U_2=12$、$U_3=-12$、$U_4=-8$（单位 V）。
再用节点电位算一遍：设下端端子为 $0$，则子图 1、2 的中间节点为 $+10$、子图 3、4 的为 $-10$；
$U_1$ 电流向上，上端比中间节点低 $2$，故 $V_{\text{上}}=8$；
$U_3$ 同法得 $V_{\text{上}}=-10-2=-12$。四个数与表一致。}

\subhdr{6.4 p38--p42：$U_{ab}=40\,\mathrm{V}$，以及「两回路独立」的严格理由}

这六页是同一张图换四套标注，最后停在 \texttt{Uab} 三个字母上，\textbf{没有答案}。补上。

\textbf{事实（从 p38--p42 的图与原文取）}：
左半是一只 $10\,\mathrm{V}$ 电压源（上端 $+$）与一只 $6\,\Omega$ 电阻\textbf{并接}在同一对母线之间；
右半是一只 $5\,\mathrm{A}$ 电流源（箭头向右）与一只 $6\,\Omega$ 电阻\textbf{并接}；
两半之间只靠\textbf{一条导线}相连（p39 把它涂红并问 「Is there any current here?」）；
$a$、$b$ 是右半那对端子的上下两端。p39--p41 三次印着同一句结论：
「The currents in the two loops are independent of each other.」

\textbf{第一步：那条红导线上的电流为零。}（这正是 §2.4 的用武之地）
把右半（$5\,\mathrm{A}$ 源与 $6\,\Omega$）整个用一个闭合面包起来，
则\textbf{穿过这个面的支路只有一条}——就是那条红导线；按广义节点 KCL，穿过闭合面的电流代数和为零，
只有一项，故该项为零。所以 p39 的问句答案是「没有电流」。

\textbf{第二步：两个回路各自闭合。}红线上无电流，左半的 $10\,\mathrm{V}$ 只能驱动左半的 $6\,\Omega$，
右半的 $5\,\mathrm{A}$ 只能在右半的 $6\,\Omega$ 里循环——这就是 p40 说的 \texttt{Two separate loops}。

\textbf{第三步：算 $U_{ab}$。}设下母线为 $0$ 电位。左半的电源直接跨在两条母线上，
故上母线为 $+10\,\mathrm{V}$。右半里 $5\,\mathrm{A}$ 全部流经它自己的 $6\,\Omega$：
电流源箭头向右，说明它把电流推向 $a$ 那一端，而 $a$ 是开路的，于是电流只能沿 $6\,\Omega$ 掉头流回左端，
即\textbf{在 $6\,\Omega$ 里由右向左}。顺着电流方向电位降低，所以 $6\,\Omega$ 右端（$a$ 一侧）比左端高
\[ 5\times6=30\,\mathrm{V} . \]
$6\,\Omega$ 左端接上母线（$+10\,\mathrm{V}$），于是
\[ V_a=10+30=40\,\mathrm{V},\qquad V_b=0\quad\Longrightarrow\quad U_{ab}=V_a-V_b=40\,\mathrm{V} . \]

\aside{p42 才把 $U_{ab}$ 的参考方向补上（由 $a$ 指向 $b$），所以按本讲义的读法 $U_{ab}=V_a-V_b$ 为正。
若把 $5\,\mathrm{A}$ 源的箭头读成向左，$U_{ab}$ 会变成 $10-30=-20\,\mathrm{V}$——
\emph{数额}由「$10$ 与 $30$ 一加一减」定，\emph{正负}由箭头定，这是两个独立的问题，
答题时最好把两件事分开写：先写 $|U_{ab}|=40\,\mathrm{V}$ 的依据，再按参考方向定号。}

\warnbox{\textbf{一句提醒}：这页不能「把左右两半等效成一个电源」来处理。
两半之间只有一条\emph{无电流}的导线，它们在电气上是\emph{共地}关系，不是串联也不是并联。
真正的等效化简要等到 1.10 等效电源定理（[C1-4]），本册只做 KVL 的直接叠加。}

%=============================================================
\hdr{§7}{1.8 支路电流法：四步，以及「独立」二字的定量含义}

\subhdr{7.1 p44--p47：目的、工具、方法，三项一样一样填}

\par\vspace{2pt}
\noindent\begin{tabular}{@{}p{2.4cm}p{1.9cm}p{11.6cm}@{}}
\tabtop
\thead{项} & \thead{页码} & \thead{课件填了什么}
\\
\tabmid
Purpose & p44--p46 & 「The purpose is to calculate the current of each branch in a complex circuit.」
（算出复杂电路里每条支路的电流）
\\
Tools & p47 & \texttt{KCL} 与 \texttt{KVL}——两个工具都已在前两节备好
\\
Methods & p47 & \texttt{Apply KCL to Nodes}（对节点用 KCL）、\texttt{Apply KVL to Loops}（对回路用 KVL）
\tabbot
\end{tabular}

\textbf{为什么需要这个方法？}一句话：1.6 的串并联化简\textbf{只能处理「能化成串并联」的电路}。
p45 那个网络里有 $R_2$ 跨接在上、中两条横支路之间，它与 $R_3$、$R_4$ 之间既不是串联也不是并联，
\textbf{化不动}。支路电流法是本课程第一个「不需要先化简」的通用方法——
它把「求电流」变成「解一个方程数正好等于未知数个数的线性方程组」。

\aside{p45 的图与 p3 的图是\textbf{同一个网络}（$I_S,R_1,R_2,R_3,R_4,U_S$ 六个元件的接法完全一致），
p3 那句「Find I by using source transformation」是 1.6 的复习题、p45 是 1.8 的入门例题。
同一个网络先用「化简」试、再用「列方程」算，正好把两种方法的适用范围比出来。
p3 的题面没有给数值，p2 的数值图接线在文字层上不能唯一复原，这两处标 [待确认]，
本讲义只给路线：$I_S$ 与 $R_1$ 是\textbf{串联}，不属于标准变换的两种组合（电压源串电阻、电流源并电阻），
所以不能直接下手；按 §7.5 的分析，只能先把 $R_1$「吸收」进电流源，然后就会在 $R_2$ 的桥式位置上停下来。}

\subhdr{7.2 p48 的四步：每一步做错会怎样}

\callout{原文（p48）：\\
(1) Directions: Specify loop direction and draw the reference direction of all currents.\\
(2) Apply KCL to the \textbf{independent} nodes.\quad
(3) Apply KVL to \textbf{independent} loops or meshes.\quad
(4) Solve the equations obtained above.\\
（原文第 (2)(3) 两步的 \emph{independent} 一词用红色突出，说明老师知道这两个字是重点。）}

\par\vspace{2pt}
\noindent\begin{tabular}{@{}p{1.0cm}p{6.6cm}p{8.4cm}@{}}
\tabtop
\thead{步} & \thead{做什么} & \thead{做错了会怎样（以及为什么不能省）}
\\
\tabmid
(1) & 在图上\textbf{标}回路绕行方向与全部支路电流的参考方向 & 省掉这一步，后面每个符号都得靠脑子记；
一条支路的方向改了而别处没改，方程组就自相矛盾。注意：方向\textbf{可以任意定}，一旦定了不许中途换
\\
(2) & 对\textbf{独立节点}列 KCL & 若把 $N$ 个节点全列上，多出来的那条方程是前几条的负和（§8.2 给实例），
方程组仍成立但白多一条；更糟的是有人会因此以为「方程不够，再补」而补错
\\
(3) & 对\textbf{独立回路或网孔}列 KVL & 若挑了两个「互相能拼成第三个」的回路，
两条方程线性相关，未知数就解不出来了。最稳的挑法是\textbf{每个网孔取一个}
\\
(4) & 解这个方程组 & 前四步都对、只有这一步错，多半是代入/消元出错；
检查办法见 §8.4 的\textbf{功率平衡}
\tabbot
\end{tabular}

\subhdr{7.3 「独立」到底独立在哪：$N-1$ 与 $M$ 从哪来}

§1.2 已经给了定理，这里补出\textbf{实用版}，因为做题时不会去算矩阵的秩：

\warnbox{\textbf{数方程的三句话}\\
（甲）节点方程：\textbf{先去掉一个节点}，其余 $N-1$ 个节点的 KCL 就是全部独立的节点方程
（去掉哪个都行，通常去掉「接地」或「连着最多支路」的那个）。\\
（乙）回路方程：\textbf{每个网孔取一条}，共 $M=B-N+1$ 条。只要每条新回路都含有前面没用过的支路，
就一定独立。\\
（丙）总数：$(N-1)+M=B$，等于未知支路电流的个数——这是「方程够用」的唯一保证。}

\aside{p52 的原句 「Only two of them are independent」 与
「We have two equations. We need three more.」 就是这条定理在 $N=3$、$B=5$ 时的特例：
$N-1=2$（三个节点方程只有两个独立）、$M=5-3+1=3$（还差三条 KVL）。
课件把这句话放在例题第 (2) 步里讲，p48 上没有写。}

\subhdr{7.4 支路里有理想电流源怎么办（p45 的图，课件没讲的那一步）}

p45 的入门例题里有一条支路是\textbf{理想电流源 $I_S$}，而 p49 起真正的例题却只有两个电压源。
课件换电路时没有交代这一步，可它恰恰是初学者最容易卡住的地方。规则两条：

\begin{enumerate}
\item \textbf{该支路的电流已知}（就是 $I_S$），所以它\textbf{不是未知数}，未知电流个数减一；
\item \textbf{该电流源两端的电压是未知的}，所以\textbf{凡是穿过它的回路，KVL 里会多出一个未知量}。
标准做法是\textbf{把 KVL 的回路选在不含这条支路的范围内}：删掉这条支路后，剩下的图仍连通，
网孔数从 $M$ 降到 $M-1$。
\end{enumerate}

拿 p45 的图数一遍：$N=4$（左、右两条竖母线各是一个节点，上横线中段一个、中横线中段一个）、
$B=6$（$I_S$、$R_1$、$R_2$、$R_3$、$U_S$、$R_4$），所以 $M=6-4+1=3$。
删掉 $I_S$ 支路后剩 5 条支路、4 个节点、仍连通，$M'=5-4+1=2$。于是
\[ \underbrace{(N-1)}_{3\ \text{条 KCL}}+\underbrace{M'}_{2\ \text{条 KVL}}=5 , \]
正好等于 5 个未知支路电流的个数（$R_1,R_2,R_3,U_S,R_4$ 五条支路的电流）。
\aside{如果硬要选穿过 $I_S$ 的回路也可以，那就得把「$I_S$ 两端电压」再设一个未知量，
方程数随之加一条——两条路等价，但前者更省事。
本册的数值例题（p55）没有电流源，所以 §8 不必处理这件事；
到了 1.9 叠加定理（§9）电流源会重新登场，那里的处理办法是「不工作时开路」。}

\subhdr{7.5 回头看 p3：1.6 的化简在这张图上能走到哪一步}

1.6 的电源等效变换只有两条规则：\textbf{电压源串电阻 $\leftrightarrow$ 电流源并电阻}。
拿这两条去套 p3/p45 那个网络（四个节点 $a,b,c,d$、六条支路），能走的路只有一条：

\begin{enumerate}
\item \textbf{第一步能走}：$I_S$ 与 $R_1$ \textbf{串在同一条上横支路}上。
一条串了电流源的支路，它的电流已经被电源钉住——$R_1$ 的存在改变不了这条支路的电流，
只能改变这条支路两端的电压。所以\textbf{从电路其它部分看，这条支路就是一个理想电流源}，
$R_1$ 可以被「吸收」进去。这一步与 §7.4 的规则同源：含理想电流源的支路，电流是\textbf{已知量}。
\item \textbf{第二步就走不动了}：$R_2$ 跨接在节点 $b$（上横线中段）与 $d$（中横线中段）之间，
它与 $R_3$（$a$--$d$）、$U_S$（$d$--$c$）、$R_4$（$a$--$c$）、$R_1$（$b$--$c$）围成一个
\textbf{桥式（非串非并）结构}。桥式结构里任何一对元件都不满足「串联」或「并联」的定义，
所以 1.6 的两条规则无处下手。
\end{enumerate}

\textbf{这就是 1.8 存在的理由}：1.6 的方法依赖「结构能一步步塌缩」，而桥式结构塌不动；
支路电流法不要求塌缩，它直接对这个结构写 $3$ 条 KCL $+$ $3$ 条 KVL（$N=4$、$B=6$、$M=3$），
把「化简」换成「解方程」。同一张图，先用 p3 的题面问「能不能化简」，再用 p45 的题面问「能不能列方程」——
课件把这张图放在这两处，用意正在于此。

\aside{p2 与 p3 是 1.6 的两道复习题，\textbf{课件没有给解}。
p3 的图与 p45 完全同源（元件、标号、甚至相对位置都一致，只是整幅图右移了约三分之一页宽），
所以本节的分析对 p3 有效；p2 的数值图（$10\,\mathrm{V}$、两只 $6\,\Omega$、$5\,\mathrm{A}$、$4\,\Omega$）
的接线在文字层上不能唯一复原，标 [待确认]，本讲义不给数值答案。}

\subhdr{7.6 p52 的三个节点标号 $a,b,c$ 为什么一次也没用上}

p52 在图上用蓝色字母标了三个节点：$a$ 在上横线左端（$U_{S1}$ 与 $R_3$ 上端交汇处）、
$b$ 在上横线右端（$R_5$ 右端与 $R_4$、$U_{S2}$ 上端交汇处）、$c$ 在下横线中部。
可是随后列出的三条方程里\textbf{只有电流、没有一个节点电位}。这不是漏写，而是方法的性质：

\begin{itemize}
\item \textbf{支路电流法的未知量是「电流」}，方程只认「方向」——哪个指针指向节点、哪个离开节点。
节点叫什么名字无关紧要，所以标号用不上；
\item \textbf{如果未知量换成「节点电位」，标号就变成主角}：以某个节点为参考（电位为零），
其余 $N-1$ 个节点的电位是未知数，KCL 写成「流出该节点的电流之和 $=0$」，
而每条支路电流用两端电位差除以电阻表示——这就是\textbf{节点电压法}，
它只需要 $N-1$ 个方程（比支路电流法的 $B$ 个少得多）。
\item 本册的课件没有讲节点电压法，p52 的 $a,b,c$ 标号只能说是一次「没用完的准备」。
但这两个字母不该浪费：记住它们，下一册或下一章讲节点电压法时，
同一个电路就会变成「以 $c$ 为参考，求 $V_a$、$V_b$」。
\end{itemize}

\aside{[教材补充] 节点电压法不在 C1-3 的内容里，此处只点一句它的方程形式，供衔接用：
以 $c$ 为参考节点（$V_c=0$）、$U_{S1}$ 的正极朝上接在 $a$ 侧，则对节点 $a$ 列「流出之和 $=0$」：
\[
\frac{V_a-V_b}{R_5}+\frac{V_a-0}{R_3}+\frac{V_a-0-U_{S1}}{R_1}=0 ,
\]
一条节点一条方程，$N-1=2$ 条就够——比五个方程省，这就是它后面被大量使用的原因。
这一句不作考点，只说明 p52 那两个字母的去向。}

%=============================================================
\hdr{§8}{主例题：五条支路的五个方程（p49--p56）}

\subhdr{8.1 图与五条方程，逐条从图上读出来}

p49 的电路：左竖支路是 $U_{S1}$（上 $+$ 下 $-$）串 $R_1$；右竖支路是 $U_{S2}$（上 $+$ 下 $-$）串 $R_2$；
中间另有两条竖支路 $R_3$、$R_4$；上横线在 $R_3$ 与 $R_4$ 之间串了 $R_5$；下横线是公共回线。
于是\textbf{两个上节点}（$R_5$ 两侧，p52 把它们标成 a 与 b）、\textbf{一个下节点}，$N=3$；
\textbf{五条支路}（左、$R_3$、$R_5$、$R_4$、右），$B=5$；$M=5-3+1=3$。

p51 标出五个电流的参考方向：$I_1$ 向上（左支路）、$I_2$ 向上（右支路）、$I_3$ 向下（$R_3$）、
$I_4$ 向下（$R_4$）、$I_5$ 向右（$R_5$）。p50 给三个网孔各画了一个\textbf{顺时针}的回环箭头。

\par\vspace{2pt}
\noindent\begin{tabular}{@{}p{1.3cm}p{2.4cm}p{1.7cm}p{4.7cm}p{4.4cm}@{}}
\tabtop
\thead{编号} & \thead{对象} & \thead{页码} & \thead{逐项（按 §2.3、§3.2）} & \thead{与课件原文对照}
\\
\tabmid
KCL 1 & 上节点 a & p52 & $I_1$ 流入，$I_3$、$I_5$ 流出 & $I_1-I_3-I_5=0$ \ 一致
\\
KCL 2 & 上节点 b & p52 & $I_5$、$I_2$ 流入，$I_4$ 流出 & $I_5+I_2-I_4=0$ \ 一致
\\
KCL 3 & 下节点 & p52 & $I_3$、$I_4$ 流入，$I_1$、$I_2$ 流出 & $I_3+I_4-I_1-I_2=0$ \ 一致
\\
KVL 1 & 左网孔 & p53 & 向上：$R_1$ 顺电流 $+R_1I_1$、$U_{S1}$ 从 $-$ 到 $+$ 取 $-U_{S1}$；
向下：$R_3$ 顺电流 $+R_3I_3$ & $I_3R_3+I_1R_1-U_{S1}=0$ \ 一致
\\
KVL 2 & 中网孔 & p53 & 向上：$R_3$ 逆电流 $-R_3I_3$；向右：$R_5$ 顺电流 $+R_5I_5$；
向下：$R_4$ 顺电流 $+R_4I_4$ & $I_5R_5+I_4R_4-I_3R_3=0$ \ 一致
\\
KVL 3 & 右网孔 & p53 & 向下：$U_{S2}$ 从 $+$ 到 $-$ 取 $+U_{S2}$、$R_2$ 逆电流 $-R_2I_2$；
向上：$R_4$ 逆电流 $-R_4I_4$ & $U_{S2}-I_2R_2-I_4R_4=0$ \ 一致
\tabbot
\end{tabular}

六条式子与 p52、p53 印的\textbf{一字不差}——这是 §3.2 那条规则的\textbf{第四次、第五次、第六次}独立通过验证
（三个网孔各一次）。三个网孔都取顺时针，所以三条 KVL 的「一致/相反」判断方式完全相同，
这也是 p50 特意把三个回环箭头画成同向的原因。

\subhdr{8.2 为什么只有两条 KCL 独立：把 p52 那句话算给他看}

把三条节点方程摆一起：
\[ I_1-I_3-I_5=0,\qquad I_5+I_2-I_4=0,\qquad I_3+I_4-I_1-I_2=0 . \]
前两条相加得 $I_1+I_2-I_3-I_4=0$，两边乘 $-1$ 就是第三条。所以\textbf{第三条是前两条的负和}，
三式只有两式独立——这正是 p52 那段粉色文字的原文：
「Only two of them are independent」 / 「We have two equations. We need three more.」
多出来的那一条\textbf{不是错的}，只是\textbf{多余}；真正的检查手段是「有没有凑够 $B$ 条」。

\subhdr{8.3 五个未知、五个方程，为什么一定有唯一解}

把 §8.1 的五个独立方程（取 KCL 1、KCL 2、KVL 1--3）写成矩阵形式，
未知向量取 $(I_1,I_2,I_3,I_4,I_5)^{\mathsf{T}}$：

\[
\begin{pmatrix}
1 & 0 & -1 & 0 & -1\\
0 & 1 & 0 & -1 & 1\\
R_1 & 0 & R_3 & 0 & 0\\
0 & 0 & -R_3 & R_4 & R_5\\
0 & -R_2 & 0 & -R_4 & 0
\end{pmatrix}
\begin{pmatrix}I_1\\I_2\\I_3\\I_4\\I_5\end{pmatrix}
=
\begin{pmatrix}0\\0\\U_{S1}\\0\\-U_{S2}\end{pmatrix}.
\]

\warnbox{\textbf{命题}\quad 当所有 $R_k>0$ 时，上式有\textbf{唯一}解。\\
\textbf{证明（换成网孔电流看，三行）}：设 $\vect{J}=(J_1,J_2,J_3)^{\mathsf{T}}$ 是三个网孔的回路电流，
令 $\vect{I}=\vect{B}^{\mathsf{T}}\vect{J}$（$\vect{B}$ 是「网孔--支路」矩阵，每个元素取 $0,\pm1$）。
KCL 的解空间维数恰为 $M=3$（§1.2），所以 $\vect{I}\leftrightarrow\vect{J}$ 是一一对应；
把它代进三条 KVL，得
\[ \bigl(\vect{B}\vect{R}\vect{B}^{\mathsf{T}}\bigr)\vect{J}=\vect{B}\vect{E}, \]
其中 $\vect{R}=\operatorname{diag}(R_1,\dots,R_5)$。对任意 $\vect{J}\ne\vect{0}$，
\[ \vect{J}^{\mathsf{T}}\vect{B}\vect{R}\vect{B}^{\mathsf{T}}\vect{J}
   =\sum_{k=1}^{5}R_kI_k^2>0 , \]
（每条支路的电阻都为正，且 $\vect{I}\ne\vect{0}$），故 $\vect{B}\vect{R}\vect{B}^{\mathsf{T}}$ 对称正定、可逆，
$\vect{J}$ 唯一，$\vect{I}$ 也唯一。}

\aside{这段证明不必背，要记住的是它的\textbf{结论与前提}：
结论是「方程数 = 未知数，且 $R_k>0$，所以有唯一解」；
前提是\textbf{线性电阻}（$U_R=RI$ 成立）。若元件是非线性的（二极管、三极管），
KCL 与 KVL 仍然成立，但方程组变成非线性方程组，唯一性不再无条件成立——
这正是课件 p4 说「Kirchhoff 定律也能用于非线性电路」的含义，也是 4.x 章要另起炉灶的原因。}

\subhdr{8.4 课件为什么停在第（4）步，以及真解出来是什么样}

p54 把四步里的最后一步标红（「Solve the equations obtained above」）就换页了，
\textbf{没有给答案}。原因不难猜：五元符号解太长，写满一屏也没人看。
本讲义用一组具体数走完第（4）步，让读者看见「解的形状」——\textbf{这组数是本讲义自拟的，课件 p49--p54 没给数值}：

\callout{[教材补充]\quad 取 $U_{S1}=12\,\mathrm{V}$、$U_{S2}=6\,\mathrm{V}$、$R_1=R_2=R_3=R_4=R_5=2\,\Omega$，
其余接法与参考方向全按 p49--p51 不变。此组数只为演示求解过程，不是课件数据。}

代入 §8.1 的六条式子（仍取 2 条 KCL + 3 条 KVL），得
\[
\begin{aligned}
I_1-I_3-I_5&=0, & I_5+I_2-I_4&=0,\\
2I_1+2I_3&=12, & 2I_5+2I_4-2I_3&=0, & 6-2I_2-2I_4&=0 .
\end{aligned}
\]
消元的顺序建议固定为：\textbf{先用 KVL 把「电源支路」的电流表出来，再用 KCL 消去中间支路}。
本题依次消去 $I_1$、$I_2$、$I_5$，最后只剩 $I_4$ 的一元一次方程，解得

\par\vspace{2pt}
\noindent\begin{tabular}{@{}p{1.3cm}p{1.7cm}p{1.3cm}p{1.7cm}p{1.3cm}p{1.7cm}p{1.9cm}p{3.6cm}@{}}
\tabtop
\thead{$I_1$} & \thead{A} & \thead{$I_2$} & \thead{A} & \thead{$I_3$} & \thead{A} & \thead{$I_4$} & \thead{$I_5=0.75\,\mathrm{A}$}
\\
\tabmid
$3.375$ & $27/8$ & $1.125$ & $9/8$ & $2.625$ & $21/8$ & $1.875$ & $15/8=0.75\,\mathrm{A}$
\tabbot
\end{tabular}

\textbf{三条校验，全部通过}（这一步不能省，它是第（4）步的唯一保险）：
（1）三条 KCL 逐条代入成立（含那条「多余」的第三条，它应给出 $0=0$）；
（2）三条 KVL 逐条代入为零；
（3）\textbf{功率平衡}：$U_{S1}$ 供出 $12\times3.375=40.5\,\mathrm{W}$，
$U_{S2}$ 供出 $6\times1.125=6.75\,\mathrm{W}$，合计 $47.25\,\mathrm{W}$；
五个电阻吸收 $2\sum I_k^2=2\times(11.39+1.27+6.89+3.52+0.56)=47.25\,\mathrm{W}$，收支相等。

\subhdr{8.5 p55：课件里唯一一组带数值的数（全解）}

p55 是全册唯一给出具体数值的题：三条竖支路并接在同一对母线之间——
左支路 $40\,\mathrm{V}$ 源（上 $+$）串 $2\,\Omega$、电流 $I_1$ 向上；
中支路 $20\,\mathrm{V}$ 源（上 $+$）串 $2\,\Omega$、电流 $I_2$ 向上；
右支路只有 $4\,\Omega$、电流 $I_3$ 向下。$N=2$、$B=3$、$M=3-2+1=2$，
所以\textbf{1 条 KCL + 2 条 KVL}就够了（三个未知电流 $I_1,I_2,I_3$）：

\par\vspace{2pt}
\noindent\begin{tabular}{@{}p{1.4cm}p{5.4cm}p{8.8cm}@{}}
\tabtop
\thead{方程} & \thead{式子} & \thead{来由}
\\
\tabmid
KCL & $I_1+I_2=I_3$ & 上节点：$I_1$、$I_2$ 流入，$I_3$ 流出
\\
KVL（外回路：左支路 $+$ 右支路） & $40-2I_1-4I_3=0$ & 顺时针：源从 $-$ 到 $+$ 取 $+40$；$R$ 顺电流取 $-2I_1$；$4\,\Omega$ 顺电流取 $-4I_3$
\\
KVL（右网孔：中支路 $+$ 右支路） & $4I_3-20+2I_2=0$ & 顺时针：$4\,\Omega$ 逆电流取 $+4I_3$；源从 $+$ 到 $-$ 取 $-20$；$R$ 逆电流取 $+2I_2$
\tabbot
\end{tabular}

（两条 KVL 用的都是「顺时针、从上母线出发」，与 p50 的习惯一致；
若有人沿逆时针走，两条式子会整体变号，解不变。）

\aside{三条支路并接在同一对母线上，只有 2 个网孔，但可以凑出 3 条回路
（左$+$中、中$+$右、左$+$右），任取两条独立的都行。本表取的是「左$+$右」与「中$+$右」，
因为这两条都穿过 $4\,\Omega$，能把 $U=4I_3$ 直接带进方程；若取「左$+$中」与「中$+$右」也完全可以，
解出来的数一样。}

解这个三元一次方程组。记两条母线之间的电压 $U=V_{\text{上}}-V_{\text{下}}$，则三条支路各给一个关系：
\[
U=40-2I_1,\qquad U=20-2I_2,\qquad U=4I_3 ,
\]
再代入 $I_1+I_2=I_3$，得
\[
\frac{40-U}{2}+\frac{20-U}{2}=\frac{U}{4}
\ \Longrightarrow\ 30-U=\frac{U}{4}
\ \Longrightarrow\ U=24\,\mathrm{V}.
\]
回代即得
\[
I_1=\frac{40-24}{2}=8\,\mathrm{A},\qquad
I_2=\frac{20-24}{2}=-2\,\mathrm{A},\qquad
I_3=\frac{24}{4}=6\,\mathrm{A}.
\]

\subhdr{$I_2=-2\,\mathrm{A}$ 这个负号在说什么}

负号只说明\textbf{实际方向与参考方向相反}：$I_2$ 的参考方向向上，实际是\textbf{向下} 2 A。
于是 $20\,\mathrm{V}$ 电源里，电流从它的 $+$ 端（上端）流入——它不是在供电，\textbf{是在被充电}，
吸收功率 $20\times2=40\,\mathrm{W}$。
而 $40\,\mathrm{V}$ 电源里电流从 $+$ 端流出，供出 $40\times8=320\,\mathrm{W}$。

\warnbox{\textbf{功率平衡校验（第（4）步的保险）}\quad
供出：$320\,\mathrm{W}$（只有 $40\,\mathrm{V}$ 那一台在供）。
吸收：$20\,\mathrm{V}$ 源 $40\,\mathrm{W}$、$2\,\Omega$（左）$2\times8^2=128\,\mathrm{W}$、
$2\,\Omega$（中）$2\times(-2)^2=8\,\mathrm{W}$、$4\,\Omega$ 支路 $4\times6^2=144\,\mathrm{W}$，
合计 $40+128+8+144=320\,\mathrm{W}$。\textbf{收支相等}，说明五个数都算对了。
另：$U=24\,\mathrm{V}$，三个支路各算一遍都得到 $24\,\mathrm{V}$（$40-16=24$、$20+4=24$、$4\times6=24$）。}

\aside{若把 p55 两台电源的极性\textbf{都}读反（$+$ 在下），整套结果\textbf{全部变号}：
$U=-24\,\mathrm{V}$、$I_1=-8\,\mathrm{A}$、$I_2=2\,\mathrm{A}$、$I_3=-6\,\mathrm{A}$，
绝对值与物理结论（谁供电、谁被充电）都不变。所以「极性读错」只影响整体符号，
不影响「做法对不对」——这也是为什么第（4）步之后一定要做一次功率平衡：
符号错了会在功率上立刻露出来（供出与吸收不再相等）。}

\subhdr{8.5.1 一个变式：把 $4\,\Omega$ 负载拿掉会怎样（两源「对顶」）}

把 p55 右支路的 $4\,\Omega$ 断开（$I_3=0$），网络里只剩两台电源与两只 $2\,\Omega$ 串联成一个回路。
这时 KCL 给出 $I_1=-I_2$，两条 KVL 合成
\[ 40-2I_1=20-2I_2=20+2I_1
\ \Longrightarrow\ 20=4I_1
\ \Longrightarrow\ I_1=5\,\mathrm{A},\ I_2=-5\,\mathrm{A}, \]
母线电压 $U=40-2\times5=30\,\mathrm{V}$。功率：$40\,\mathrm{V}$ 源供出 $40\times5=200\,\mathrm{W}$，
$20\,\mathrm{V}$ 源吸收 $20\times5=100\,\mathrm{W}$（电流从它的 $+$ 端流入），
两只 $2\,\Omega$ 各吸收 $2\times5^2=50\,\mathrm{W}$，合计 $100+50+50=200\,\mathrm{W}$，收支相等。

\warnbox{\textbf{这个变式说明一件常被误解的事}：\textbf{没有负载，一样可以有电流}。
两台电压源并接在同一对母线上，只要电压不相等，高的一台就会通过两个内阻给低的一台「充电」，
哪怕外面什么都没接。p55 里 $I_2=-2\,\mathrm{A}$ 是这同一件事的弱化版（有负载分流）；
把负载拿掉，它就变成 $I_2=-5\,\mathrm{A}$，一望而知。所以\textbf{看到负电流先想「谁在给谁充电」，
而不是「我算错了」}。}



\subhdr{8.6 p56：六电流那张图，只做方程数讨论}

p56 是全册最后一个例题图，五种电阻（$4\,\Omega$、$5\,\Omega$、$2\,\Omega$、$3\,\Omega$、$1\,\Omega$）、
两台电源（$5\,\mathrm{V}$、$10\,\mathrm{V}$，后者左 $+$ 右 $-$）、六个电流标号（$I_1\sim I_6$），
\textbf{没有文字、没有答案}。本讲义不硬解它，原因要说清：

\gapitem{[待确认] 两条读法不能同时成立}{
按扫描图的连线看，中横线的两端与左右两条竖边相交，这会把上、中、下三层\textbf{短接}；
可图上又把 $I_4$、$I_5$ 标在左右竖边上，说明这两条竖边\textbf{本身应当是有元件的支路}。
两种读法互相冲突，原图的分辨率下无法判定（文字层只留下元件数值与电流标号的位置，没有连线）。
因此\textbf{本页只做方程数的讨论，不做数值求解}。}

\textbf{方程数怎么数（这一点与拓扑读法无关，永远成立）}：先数 $B$，再数 $N$，用 $(N-1)+M$ 检查，
其中 $M=B-N+1$。按「左右竖边各算一条支路」的读法：$B=8$（$4\,\Omega$、$5\,\Omega$、$5\,\mathrm{V}$ 支路、
$2\,\Omega$、$3\,\Omega$、左竖边、右竖边、底部 $1\,\Omega+10\,\mathrm{V}$），$N=6$，
于是 $M=8-6+1=3$，需要 $(6-1)+3=8$ 条方程——而图上只标了 6 个电流，
\textbf{缺 $4\,\Omega$ 与 $5\,\Omega$ 两条支路的电流}。若按「三层母线短接」的读法则 $B=6$、$N=4$、$M=3$，
需要 6 条方程，与 6 个标号对得上，但那与 $I_4$、$I_5$ 的存在矛盾。
两种读法都记在这里，供对照原图核对。

\subhdr{8.7 方法选择：这张表决定「先化简还是先列方程」}

本册先后出现了三种手段，各有各的地盘。拿到一道题，先按这张表定位：

\par\vspace{2pt}
\noindent\begin{tabular}{@{}p{5.2cm}p{3.6cm}p{6.0cm}@{}}
\tabtop
\thead{电路的样子} & \thead{首选手段} & \thead{为什么}
\\
\tabmid
纯串并联、一眼能认出哪些是串、哪些是并 & 1.6 的等效化简 & 结构能一步步塌缩到只剩一个电阻
\\
含跨接支路（桥式），塌不动 & \textbf{支路电流法}（本册 §7--§8） & 不需要塌缩，直接解 $B$ 元一次方程组
\\
含理想电流源 & 支路电流法 $+$ §7.4 两条规则 & 该支路电流已知，KVL 回路避开它
\\
线性、多台独立源、只求某个量 & 1.9 叠加定理（§9） & 每个源单独作用，电路变得更简单
\\
只求某一支路／只求端口电压 & 1.10 等效电源定理（[C1-4]） & 把「其余部分」压成一个电源加一个电阻
\tabbot
\end{tabular}

\warnbox{\textbf{选择之前先读题}：题目要的是「\emph{全部}支路电流」还是「\emph{某一个}量」？
要全部——支路电流法最省事，$B$ 个未知、$B$ 个方程，一次解完；
只要一个——用等效化简或等效电源定理，别把整张网的电流都解出来。
p37 就是后者的例子：题目只要 $U_{AB}$，图上又给了那条支路的电流，
于是连右边的 $2\,\Omega/20\,\mathrm{V}$ 与 $4\,\Omega$ 两支路都不必碰（§5.4）。}

%=============================================================
\hdr{§9}{1.9 叠加定理预览（p57--p67）与全册公式卡、易错点}

课件在 p56 之后插了一张分节页（p57 只有 「Preview / 1.9 Superposition theorem」），
然后用 11 页把 1.9 的「定义—内容—步骤—注意」四块各讲了一遍。
这 11 页与下一册 [C1-4] 的前 10 页内容重叠，但按「课件出现过的内容都要有落点」的规矩，
本节把这 11 页逐页交代清楚；下一册讲例题时再深挖。

\subhdr{9.1 定义（p58）：一句话里藏了三个条件}

\callout{原文（p58）：\quad「The voltage across (or current through) an element in a linear electric
circuit is the algebraic sum of the voltages across (or currents through) that element due to
each independent source acting alone.」\\
译：\textbf{线性电路中任一元件上的电压（或电流），等于各独立电源\emph{单独作用}时在该元件上产生的
电压（或电流）的\emph{代数和}}。\\
同页还印了一句解释与一句目的：「It means that each independent source has its effect on each circuit passive elements. And the final effect can be found using the algebraic sum of all the above effects.」 与 「Purpose: Simplify the calculation」（目的：简化计算）。}

三个条件必须一个个点出来，因为用错任意一个都会得出错答案：
\begin{enumerate}
\item \textbf{linear（线性）}：电路必须是线性的。所以二极管、三极管那种「$U$ 与 $I$ 不成比例」的元件不能直接套；
\item \textbf{algebraic（代数）}：叠加时\textbf{带正负号相加}，不是把绝对值加起来（p67 第 (3) 条说的「注意方向」正是这个意思）；
\item \textbf{each ... acting alone}：每次只让\textbf{一个}独立源工作，其余全部「不工作」（怎么算不工作，见 §9.2）。
\end{enumerate}

\subhdr{9.2 内容（p59--p60）：从「定义为」到「怎么做」的那一步}

\callout{原文（p59，全段）：\quad「If a number of voltage or current sources are acting in a linear network,
the resulting current in any branch is the algebraic sum of all the currents that would be produced in it,
when each source acts alone, \textbf{all the other independent sources are replaced by their internal
resistances.}」\\
同页图下另印一行：$I=I'+I''+I'''+I''''$（四个撇号对应四台独立源：三台电压源 $U_{S1},U_{S2},U_{S3}$
加一台电流源 $I_S$，与 p59 图上数出来的源数一致）。}

\textbf{这一段的裂缝在这里}：它说其余源「\emph{用内阻代替}」，可内阻是多少、怎么在图上画出来，一句没说。
p60 把这句话用黄色高亮又追问了一遍（原文照录）：
「What do you mean by saying all the other independent sources are replaced by their internal resistances? How to deal with the other sources when they are not working?」
p61 才给出答案——而答案与「内阻」二字之间还差一步换算：

\warnbox{\textbf{电源「不工作」的等效}\quad
\textbf{理想电压源不工作 $=$ 短路}（因为理想电压源的内阻为 $0$，用 $0\,\Omega$ 代替它，就是一根导线）；\\
\textbf{理想电流源不工作 $=$ 开路}（因为理想电流源的内阻为 $\infty$，用断路代替它，就是把它拔掉）。\\
课件 p61 的两句原文：「For ideal voltage source, it is like a short circuit when it is not working.」
与 「For ideal current source, it is like a open circuit when it is not working.」}

\aside{「用内阻代替」与「短路/开路」的关系就是取极限：
实际电压源 $=$ 理想电压源串内阻 $r$，$r\to0$ 时「串一个 $0\,\Omega$」等于什么都没串；
实际电流源 $=$ 理想电流源并内阻 $r$，$r\to\infty$ 时「并一个无穷大电阻」等于断开。
所以两句话是同一件事的理想化版本，课件只写了结论、没写这一步换算。}

\subhdr{9.3 p62--p65 的四张图：同一个操作的四个镜头}

这四页的图几乎一样，差别只在「哪些源的符号还在」，看懂了它们其实是一套连续动作：

\par\vspace{2pt}
\noindent\begin{tabular}{@{}p{1.2cm}p{5.9cm}p{8.5cm}@{}}
\tabtop
\thead{页} & \thead{图上发生了什么} & \thead{它在演示第几步}
\\
\tabmid
p61 & 完整原图（三台电压源、一台电流源都在），左上加框 「Assume only this one is working」 指住 $U_{S1}$ & 起手式：选定「只有一个在工作」的那一个
\\
p62 & $U_{S2}$、$U_{S3}$ 的圆圈消失、只剩 $+/-$ 与标号（\textbf{被短路}）；$I_S$ 处只剩一个断口（\textbf{被开路}） & 把「不工作」落到图上
\\
p63 & 连 $U_{S2}$、$U_{S3}$ 的标号也去掉，右半只剩 $R_4$ 与 $R_2$ 两条电阻 & 化简到「能算」的样子
\\
p64 & 三个电源符号全部画回，回到原图 & 复原，准备换下一个源
\\
p65 & 换成 $I_S$ 被红圈标住，电压源只剩极性记号 & 同一个动作换主角再做一遍
\tabbot
\end{tabular}

\subhdr{9.4 步骤（p66）与四点注意（p67）}

\callout{原文（p66）：\\
(1) Turn off all the other independent sources except the one in operation.
Find the outputs (voltage or current) due to that active source.\\
(2) Repeat step (1) for each of the other independent source.\\
(3) Find the total contribution by adding algebraically all the contributions due to the independent sources.}

三步与 §7.2 的支路电流法四步是两种不同的「拆」法：支路电流法\textbf{拆方程}（一个电路、多个方程），
叠加定理\textbf{拆电路}（一个电路、多个更简单的电路）。前者方程多但每个都简单，后者电路多但每个只用算一遍串并联。

p67 的四点注意（原文照录后逐条解读）：

\par\vspace{2pt}
\noindent\begin{tabular}{@{}p{0.8cm}p{7.2cm}p{7.6cm}@{}}
\tabtop
\thead{条} & \thead{课件原文（节录）} & \thead{本讲义的解读}
\\
\tabmid
(1) & 「It is for lineal circuit. In order to use the superposition theorem with circuit currents and voltages, all of the components must be linear」 & 线性是前提。
「lineal」是「linear」的拼写错误，原文如此
\\
(2) & 「It is not applicable to power. Since power is not a linear quantity.」 &
\textbf{功率不能叠加}。原因见下方的最小反例
\\
(3) & 「Pay attention to the direction. It is not applicable to power, since power is not a linear quantity. ...」 & 前半句对（叠加要带符号）；\textbf{后半句是第 (2) 条的原文照抄}，
与「方向」无关，属排版时的复制错误
\\
(4) & 「Each group can have more than one sources.」 & 可以\textbf{成组}处理：
把几台源一起保留、一起关掉，只要每次只保留一组
\tabbot
\end{tabular}

\warnbox{\textbf{为什么功率不能叠加：一个最小反例（[教材补充]，一两行就够）}\quad
取 $1\,\Omega$ 电阻，两台各给 $1\,\mathrm{A}$ 的电流源。第一个源单独作用：$I=1\,\mathrm{A}$，
功耗 $1^2\times1=1\,\mathrm{W}$；第二个源单独作用也是 $1\,\mathrm{W}$；两者相加得 $2\,\mathrm{W}$。
但两个源同时作用时电流是 $2\,\mathrm{A}$，真实功耗 $2^2\times1=4\,\mathrm{W}\ne2\,\mathrm{W}$。
差别来自 $p=Ri^2$ 是\textbf{二次}关系：$(i_1+i_2)^2\ne i_1^2+i_2^2$。
所以叠加定理只对\textbf{一次量}（电压、电流）成立，对功率必须先叠加电流（或电压）再平方。}

\subhdr{9.5 公式卡（按做题顺序排，不是按课件顺序）}

\par\vspace{2pt}
\noindent\begin{tabular}{@{}p{2.2cm}p{7.0cm}p{2.0cm}p{4.4cm}@{}}
\tabtop
\thead{量／判据} & \thead{式子} & \thead{依据} & \thead{前提}
\\
\tabmid
网孔数 & $M=B-N+1$ & p7--p10、p49 & 连通平面电路
\\
方程总数 & $(N-1)+M=B$ & §1.2（本讲义补） & 每条支路一个未知电流
\\
KCL & $\sum I_{\text{入}}=\sum I_{\text{出}}$ & p13、p14 & 集中参数
\\
广义节点 & 穿过任一闭合面的电流代数和为零 & p15--p17 & 闭合面内不累积电荷
\\
KVL & $\sum(\pm U_k)=0$，同向取 $+$ & p20、p23 & 电位单值
\\
电阻项符号 & 顺绕行方向上的电流取 $+R_kI_k$，逆取 $-R_kI_k$ & p32、p35（本讲义补规则） & 参考方向已标注
\\
电源项符号 & 绕行从 $+$ 到 $-$ 取 $+U_S$，反之取 $-$ & p23、p27 & 极性已标注
\\
端口电压 & $U=U_S+RI$（电流与 $U$ 同向时） & p32 & 单条电源--电阻支路
\\
端口电压 & $U=U_S-RI$（电流反向时） & p35 & 同上
\\
唯一解 & $\det\ne0$，因 $\vect{B}\vect{R}\vect{B}^{\mathsf{T}}$ 正定 & §8.3（本讲义补） & 所有 $R_k>0$
\\
不工作的电压源 & 短路 & p61 & 理想源
\\
不工作的电流源 & 开路 & p61 & 理想源
\\
功率能否叠加 & \textbf{不能}，先叠 $I$ 再算 $RI^2$ & p67 (2) & 任何情况
\tabbot
\end{tabular}

\subhdr{9.6 本册全部符号判断的汇总：七处验证，同一条规则}

本讲义 §3.2 的规则在七个地方各被判过一次，全部对上。这张表可以用来快速回忆「哪个页对应哪种情形」：

\par\vspace{2pt}
\noindent\begin{tabular}{@{}p{1.6cm}p{5.6cm}p{4.4cm}p{4.0cm}@{}}
\tabtop
\thead{页码} & \thead{情形} & \thead{规则给出的结果} & \thead{课件的式子}
\\
\tabmid
p23 & 四电阻回路，顺时针绕行 & $U_1-U_2+U_3-U_4=0$ & 与课件一字不差
\\
p27 & 两电源两电阻单回路 & $U_{S2}+R_2I_2+R_1I_1-U_{S1}=0$ & 一致（末项下标笔误）
\\
p32 & 端口支路，电流向下 & $U_S+RI-U=0$ & 一致
\\
p35 & 端口支路，电流向上 & $U_S-RI-U=0$ & 一致
\\
p36 & 四个子图（极性 $\times$ 电流方向） & $U_1=U_S-RI$、$U_2=U_S+RI$ 等四式 & 课件只给四式、未给对应
\\
p53 & 三网孔同向绕行 & 三条 KVL（§8.1） & 三条全一致
\\
p55 & 两源并联供一负载 & $I_1=8$、$I_2=-2$、$I_3=6\,\mathrm{A}$ & 课件未给答案
\tabbot
\end{tabular}

\subhdr{9.7 与下一册的接口：这一册留下的三个线头}

本册在 1.9 处收尾，但有三处「只起了头、没接完」的地方，下一册 [C1-4] 直接从这里接着走：

\begin{enumerate}
\item \textbf{p28 的端口电压 $U$}：它是 1.10 等效电源定理里「开路电压」的雏形。
本册只做了「$U$ 与 $U_S$、$RI$ 的关系」，没有再问「这个 $U$ 能不能用一个电源代替」。
\item \textbf{p59 的 $I=I'+I''+I'''+I''''$}：四个撇号对应四台源，本册只交代了「每台单独作用怎么做」，
没有做例题；例题在 [C1-4] 的 E1--E5。
\item \textbf{p61 的短路/开路规则}：它是叠加定理的操作核心，也是 1.10 里「把独立源置零求等效电阻」的同一条规则。
\end{enumerate}

\aside{[教材补充] 一句话把 1.10 的轮廓说清（C1-3 未涉及，只为让读者知道这条线通向哪里）：
一个含源线性二端网络，对外可以等效成「一个电压源串一个电阻」，其中电压源取\textbf{开路电压}
$U_{OC}$、电阻取该网络内部独立源全部置零后的等效电阻 $R_{eq}$。本册 §5.3 求的就是 $U_{OC}$ 的特例，
§9.2 的「短路/开路」就是「置零」的操作；剩下的 $R_{eq}$ 在 [C1-4] 讲。这一句不作考点。}

\subhdr{9.8 本册英文术语对照（读课件时查这一张）}

课件全英文、教材中文，最容易卡住的不是公式而是词。下表把本册出现的术语按出现页码列齐，
右边的「注意」一列写的是容易译错或容易混的地方：

\par\vspace{2pt}
\noindent\begin{tabular}{@{}p{3.9cm}p{2.6cm}p{3.2cm}p{5.0cm}@{}}
\tabtop
\thead{英文（页码）} & \thead{中文} & \thead{出现处} & \thead{注意}
\\
\tabmid
node / branch & 节点 / 支路 & p6 & 节点要「三条及以上」支路相交（§1.1）
\\
loop / mesh & 回路 / 网孔 & p6 & 网孔是「内部无支路」的回路，数目唯一
\\
independent node / loop & 独立节点 / 独立回路 & p48 & 「独立」指线性无关，不是「单独一个」
\\
closed boundary & 闭合边界（闭合面） & p15--p17 & 广义节点的边界，不是「闭合回路」
\\
generalized node & 广义节点 & p11、p15 & 把一整块电路当一个节点
\\
generalized loop & 广义回路 & p28 & 把「端口电压」那一段当回路的一条边
\\
imaginary loop & 假想回路 & p29 & 与广义回路是同一件事的两种叫法
\\
loop direction & 绕行方向 & p19、p22 & 与「电流参考方向」不是一回事（§3.2）
\\
reference direction & 参考方向 & p48、p12 & 图上箭头，可任意定，算出负号不算错
\\
algebraic sum & 代数和 & p13、p20 & 带正负号相加，不是绝对值相加
\\
KCL / KVL & 基尔霍夫电流/电压定律 & p11、p18 & 前者管节点，后者管回路
\\
\tabbot
\end{tabular}

{\small 下表接着上一半（后半是 1.8 与 1.9 的术语）：}

\par\vspace{2pt}
\noindent{\small\begin{tabular}{@{}p{3.9cm}p{2.5cm}p{3.0cm}p{4.6cm}@{}}
\tabtop
\thead{英文（页码）} & \thead{中文} & \thead{出现处} & \thead{注意}
\\
\tabmid
BCA (branch current analysis) & 支路电流法 & p44 & 本册 1.8 的主角
\\
source transformation & 电源等效变换 & p1--p3 & 1.6 的内容，只对「源串电阻/源并电阻」成立
\\
practical voltage / current source & 实际电压源 / 实际电流源 & p1 & 理想源串/并内阻
\\
ideal (voltage/current) source & 理想电压源 / 理想电流源 & p61 & 内阻 $0$ / $\infty$
\\
internal resistance & 内阻 & p60 & 「用内阻代替」$=$ 短路/开路（§9.2）
\\
short circuit / open circuit & 短路 / 开路 & p61 & 电压源不工作短路，电流源不工作开路
\\
superposition theorem & 叠加定理 & p57 & 只对一次量成立，功率不行
\\
independent source & 独立源 & p58 & 与受控源相对；本册只出现独立源
\\
linear (circuit) & 线性（电路） & p58、p67 & 原文 p67 误拼为 lineal
\\
terminal / port & 端子 / 端口 & p28、p37 & 端口电压 $U$ 的两个端点
\\
bridge & 桥式（结构） & 隐含于 p3、p45 & 非串非并，化简法在此失效（§7.5）
\tabbot
\end{tabular}}

\subhdr{9.9 p58--p59 那两张示例电路：本册只读图，例题留给下一册}

p58 一页放了\textbf{两张}带数值的电路（左下与右下），p59 又画了一张四支路网络，
它们都是 1.9 例题的题面。本册只把「读图要读什么」交代清楚，计算留给 [C1-4] 的 E1--E5
（下一册的 p12 起有 E1、E2，p19 起有 E3--E5，与本册这三张图一一对应）：

\par\vspace{2pt}
\noindent\begin{tabular}{@{}p{4.3cm}p{6.0cm}p{5.4cm}@{}}
\tabtop
\thead{图} & \thead{图上给了什么} & \thead{读图时盯住哪两处}
\\
\tabmid
p58 左下 & $7\,\mathrm{A}$ 电流源、$2\,\Omega$、$3\,\Omega$、$1\,\Omega$、$35\,\mathrm{V}$ 电压源（左 $+$ 右 $-$）、
$4\,\Omega$，待求电流 $I$ 标在 $4\,\Omega$ 上、方向向右 & 每台源的参考方向（$7\,\mathrm{A}$ 斜向左下、
$35\,\mathrm{V}$ 左正）；$I$ 的箭头在哪条支路
\\
p58 右下 & $3\,\mathrm{A}$ 与 $5\,\mathrm{A}$ 两台电流源（都向上）、$4\,\Omega$、$2\,\Omega$、$3\,\Omega$、
$10\,\mathrm{V}$（上 $+$）、待求 $I$ 标在 $10\,\mathrm{V}$ 支路左侧、方向向左 & 两台电流源\textbf{同向}
（都是往上注）；$I$ 的箭头方向与 $10\,\mathrm{V}$ 的极性
\\
p59 & 三条支路各含一台电压源（$U_{S1},U_{S2},U_{S3}$，全为上 $+$ 下 $-$）串一只电阻（$R_1,R_3,R_2$），
上横线串一台电流源 $I_S$，待求 $I$ 在 $R_3$ 支路、方向向下 & $I=I'+I''+I'''+I''''$ 里四个撇号
各对应哪一台源；$I_S$ 串在横线上意味着它\textbf{把网络分成左右两块}
\tabbot
\end{tabular}

\aside{这三张图的共同点：\textbf{每台源都单独标了参考方向，待求量只有一条支路上的一个电流}。
这正是叠加定理最适合的场合——如果题目要「全部支路电流」，用支路电流法（§8）反而更省。
选择依据见 §8.7 的表。}

%=============================================================
\hdr{§10}{缺口清单：课件哪句话跳了步，本讲义补在哪}

下面 18 条按「类型｜页码｜课件原文短语｜缺了什么｜补在哪」写，每条都能回到某一页的某一句话。
类型的分布是：机制缺口 5 条、拼接缺口 3 条、答案形态 4 条、形式改写 1 条、适用前提 2 条、印刷笔误 2 条、版面重复 1 条。

\subhdr{机制缺口}

\gapitem{机制缺口}{p6\quad 「Mesh: Independent loop.」\quad
缺：\textbf{「独立」独立在哪}、网孔有几个、$M$ 与 $N,B$ 的关系\quad
补：§1.2 给出定理（$N-1$ 条 KCL $+M$ 条 KVL $=B$）并给机制说明；§1.3 用 p7--p10 逐页核对}

\gapitem{机制缺口}{p11--p12\quad 五项提纲（For / Content / Examples / Other expression /
Generalized node）下面全是空的\quad
缺：KCL \textbf{为什么成立}\quad
补：§2.2 用「电荷守恒 $+$ 集中参数，节点无处存电荷」给出理由}

\gapitem{机制缺口}{p19--p23\quad \texttt{Loop direction} 与 $U_1-U_2+U_3-U_4=0$\quad
缺：\textbf{元件电压的正负怎么定}——课件把箭头画好、把答案印在下面，规则一句没写\quad
补：§3.2 给出统一规则（元件箭头与绕行方向一致取正、相反取负），§3.3 逐段验证}

\gapitem{机制缺口}{p47--p48\quad 「Apply KCL to the \emph{independent} nodes」 /
「Apply KVL to \emph{independent} loops or meshes」（independent 特意标红）\quad
缺：什么是独立、独立节点数与独立回路数各几个\quad
补：§7.3 给出「去掉一个节点 + 每个网孔取一条」的实用数法；§8.2 用 p52 的三式演示为什么会多一条}

\gapitem{机制缺口}{p61\quad 「For ideal voltage source, it is like a short circuit when it is not working.」\quad
缺：为什么「用内阻代替」等价于「短路/开路」\quad
补：§9.2 旁注给出 $r\to0$、$r\to\infty$ 的换算}

\subhdr{拼接缺口}

\gapitem{拼接缺口}{p13\quad 「the algebraic sum of the current entering a node」\quad
缺：\textbf{符号规则}。p13 说「代数和」，p14 才用一个例子暗示，中间没有一句话说「流入取正还是流出取正」\quad
补：§2.3 把三种等价写法列出并对照 p14 逐项核验}

\gapitem{拼接缺口}{p45 与 p49\quad p45 印着「The purpose is to calculate the current of each branch
in a complex circuit.」并配一张\emph{含理想电流源}的图；p49 的正片例题却只有两个电压源\quad
缺：\textbf{换成只有电压源的电路为什么合理}、有电流源时怎么列方程\quad
补：§7.4 给出两条规则（电流已知 $\to$ 不当未知数；电流源端电压未知 $\to$ KVL 回路避开它），并数出 $3+2=5$}

\gapitem{拼接缺口}{p50--p54\quad 例题只讲完四步，「Solve the equations obtained above」 之后直接换页\quad
缺：\textbf{这五个方程怎么解}、解出来什么形状\quad
补：§8.3 给矩阵与唯一性证明，§8.4 用一组自拟数（[教材补充]）走完全过程并做功率校验}

\subhdr{答案形态}

\gapitem{答案形态}{p2--p3\quad 「Find I by using source transformation.」\quad
缺：两道 1.6 复习题\textbf{只出题不给解}；也没说 p3 与 p45 是同一个网络\quad
补：§7.5 说明这张图在 1.6 只能化简一步（$I_S$ 串 $R_1$ 可并成理想电流源）、
到桥式结构处就走不通，并据此解释 1.8 为什么要换方法；p2 的数值图标 [待确认]}

\gapitem{答案形态}{p36\quad \texttt{Find U1，U2，U3，U4} 与四行「$U\ \square=U_S+IR$」\quad
缺：$U_1\sim U_4$ 与四个式子的\textbf{对应关系}（方框留空）\quad
补：§6.3 给出对应表（$U_1\leftrightarrow(3)$、$U_2\leftrightarrow(1)$、$U_3\leftrightarrow(2)$、
$U_4\leftrightarrow(4)$）与一组数值自检}

\gapitem{答案形态}{p37--p42\quad \texttt{Find UAB}、\texttt{Find Uab}、「Is there any current here?」\quad
缺：六页\textbf{一个答数都没有}；p39 的问句也没有回答\quad
补：§5.4 给出 $U_{AB}=12\,\mathrm{V}$（沿已知电流的支路走一遍）；§6.4 给出 $U_{ab}=40\,\mathrm{V}$，
并用广义节点 KCL 证明那条红导线电流为零}

\gapitem{答案形态}{p55--p56\quad 两页只有元件数值与电流标号（p55 的「40 V / 20 V / 2 $\Omega$ / 2 $\Omega$ / 4 $\Omega$」、
p56 的「4 $\Omega$ / 5 $\Omega$ / 5 V / 2 $\Omega$ / 3 $\Omega$ / 1 $\Omega$ / 10 V」），没有一个动词\quad
缺：这两张图要算什么\quad
补：§8.5 把 p55 全解（$I_1=8$、$I_2=-2$、$I_3=6\,\mathrm{A}$，$U=24\,\mathrm{V}$，功率收支 320 W）；
p56 因拓扑读法冲突标 [待确认]，只做方程数讨论（§8.6）}

\subhdr{形式改写}

\gapitem{形式改写}{p6\quad 「Mesh: Independent loop.」\quad
缺：这句是\textbf{常用形式}（「网孔 $=$ 平面电路里内部无支路的回路」）的压缩，
原形式应是「独立回路数 $=B-(N-1)$」这个\emph{关系}，而不是「网孔」这个\emph{对象}\quad
补：§1.2 先写原形式（$N-1$ 与 $B-(N-1)$），再推出常用形式 $M=B-N+1$，并标明课件用的是常用形式}

\subhdr{适用前提}

\gapitem{适用前提}{p9\quad 「N: 4\quad B: 6\quad L: 7\quad M: 3」\quad
缺：$L=7$ 的计数规则（「简单回路」怎么数）——四页答案只给了数，没给数法；
而 $L$ 依赖桥臂怎么连，恰恰是最难数的一项\quad
补：§1.3 用 p7、p8、p10 三页把「$M$ 有公式、$L$ 只能数」讲清，p9 的 $L=7$ 标 [待确认] 不复核}

\gapitem{适用前提}{p20、p56\quad 「algebraic sum of all voltages around a closed path (or loop) is zero」
与 p56 那张六电流图\quad
缺：KVL 的前提（电位单值、集中参数）与 $M$ 公式的前提（连通、平面）；
p55--p56 这类「多条含源支路并联」的题还要求所有 $R_k>0$ 才有唯一解\quad
补：§3.4、§1.2、§8.3 分别点明前提；§9.5 公式卡最后一列专门列「前提」}

\subhdr{印刷与记号}

\gapitem{印刷笔误}{p27\quad \texttt{R1I1－U1S＝0}\quad
缺：末项下标应为 $U_{S1}$（同页上方标号是 $U_{S1}$）\quad
补：§4.3 点名，并说明它不影响方程结构}

\gapitem{印刷笔误}{p67\quad 「(1) It is for lineal circuit」 与 「(3) Pay attention to the direction. It is not applicable to power, since power is not a linear quantity ...」\quad
缺：(1) 的 \texttt{lineal} 应为 \texttt{linear}；第 (3) 条的说明句是第 (2) 条的\textbf{原文照抄}，
与「方向」无关\quad
补：§9.4 逐条解读时点名，并把第 (3) 条该说的话补上（叠加是代数和，方向相反要相减）}

\subhdr{版面}

\gapitem{版面重复}{p1 与 p43、p3 与 p45、p45 与 p46、p61 与 p64\quad 同一张图重复出现\quad
缺：内容只应计一次；p3 与 p45 同一网络分别挂在 1.6 复习与 1.8 入门下，容易误以为是两个题\quad
补：§0.2 备案；§7.1 旁注说明这两页是同一个网络的两个入口；§0.1 的坐标表按去重后的内容计}

%=============================================================
\hdr{§11}{三句话总结}

\begin{enumerate}
\item \textbf{两句话一条规则}：KCL 说「节点上电流的代数和为零」，KVL 说「回路上电压的代数和为零」；
「代数和」三个字的落地形式只有一条——\textbf{箭头与绕行方向（或与「流入/流出」）一致取正、相反取负}，
其余全是记账（§2.3、§3.2）。
\item \textbf{方程数是被拓扑钉死的}：$N$ 个节点只给 $N-1$ 条独立 KCL，$B$ 条支路给出 $M=B-N+1$ 个网孔，
合起来恰好 $B$ 条——等于未知支路电流的个数；所以支路电流法\textbf{永远不多不少}，
多出来的一定是冗余、少一条一定是漏了（§1.2、§7.3、§8.2）。
\item \textbf{算完必须验}：KCL、KVL 各回代一遍，再做一次\textbf{功率平衡}
（供出 $=$ 吸收）；p55 的 $320=320$ 与本讲义 §8.4 的 $47.25=47.25$ 都是这么来的。
负号不是错，它只是说「实际方向与箭头相反」（§8.5）。
\end{enumerate}

%=============================================================
\hdr{§12}{与课件的页码对照表}

\par\vspace{2pt}
\noindent\begin{tabular}{@{}p{1.9cm}p{6.6cm}p{7.4cm}@{}}
\tabtop
\thead{课件页} & \thead{课件内容} & \thead{本讲义位置}
\\
\tabmid
p1 & Review+Preview 目录页（与 p43 同一张） & §0.1、§0.2
\\
p2--p3 & 1.6 电源等效变换复习题（只出题） & §0.2、§7.1 旁注（标 [待确认]）
\\
p4--p5 & Kirchhoff 与 Bunsen 生平、光谱分析 & §2.2 旁注一处（不复述生平）
\\
p6 & Node / Branch / Loop / Mesh 四个定义 & §1.1、§1.2
\\
p7--p10 & 四页计数练习（N、B、L、M） & §1.3
\\
p11 & KCL 五项提纲（空） & §2.1
\\
p12 & 节点电流图的五个箭头 & §2.3
\\
p13 & KCL 陈述（$\Sigma i=0$、$\Sigma I=0$） & §2.1、§2.2
\\
p14 & 例：$I_1+I_3-I_2-I_4-I_5=0$ 与移项写法 & §2.3
\\
p15--p17 & 广义节点（closed boundary）：三支、三角、四边 & §2.4
\\
p18 & KVL 的对象：Loop / Mesh & §3.1
\\
p19 & Loop direction（绕行方向） & §3.1、§3.2
\\
p20 & KVL 陈述（$\Sigma u=0$、$\Sigma U=0$） & §3.1
\\
p21--p23 & E1：四电阻标极性 $\to$ 加绕行 $\to$ $U_1-U_2+U_3-U_4=0$ & §3.3、§3.4
\\
p24--p27 & E2：两电源两电阻单回路 $\to$ $U_{S2}+R_2I_2+R_1I_1-U_{S1}=0$ & §4
\\
p28--p29 & Imaginary loop、端口电压 $U$ & §5.1
\\
p30--p31 & 定方向：绕行、电流 $I$、电压参考方向 & §5.2
\\
p32 & $U_S+RI-U=0$ $\to$ $U=U_S+RI$ & §5.3
\\
p33--p35 & 电流反向 $\to$ $U_S-RI-U=0$ $\to$ $U=U_S-RI$ & §6.1、§6.2
\\
p36 & Find $U_1\sim U_4$（四小题 + 四行待填式） & §6.3
\\
p37 & Find $U_{AB}$（$2\,\Omega$、$14\,\mathrm{A}$、$40\,\mathrm{V}$） & §5.4
\\
p38--p42 & Find $U_{ab}$（$10\,\mathrm{V}$、$5\,\mathrm{A}$、两只 $6\,\Omega$） & §6.4
\tabbot
\end{tabular}

\subhdr{续表：p43 起的 1.8 与 1.9}

\par\vspace{2pt}
\noindent\begin{tabular}{@{}p{1.9cm}p{6.6cm}p{7.4cm}@{}}
\tabtop
\thead{课件页} & \thead{课件内容} & \thead{本讲义位置}
\\
\tabmid
p43 & Review+Preview 目录页（重复） & §0.2
\\
p44--p47 & 1.8.1：目的、工具（KCL+KVL）、方法 & §7.1
\\
p48 & 1.8.2：四步 & §7.2、§7.3
\\
p49 & 1.8.3：数节点/支路/回路/网孔与「5 equations」 & §1.4、§8.1
\\
p50 & 三个网孔的绕行方向 & §8.1
\\
p51 & 五个电流的参考方向 & §8.1
\\
p52 & 三条 KCL + 「只有两条独立」 & §8.2
\\
p53--p54 & 三条 KVL、合成五式、第（4）步 & §8.1、§8.3、§8.4
\\
p55 & 两电源并联供 $4\,\Omega$（无数值答案） & §8.5
\\
p56 & 六电流多网孔图（无要求） & §8.6（标 [待确认]）
\\
p57 & 1.9 分节页 & §9 开头
\\
p58 & 1.9.1 定义与目的 & §9.1
\\
p59--p60 & 1.9.2 内容、$I=I'+I''+I'''+I''''$、两个追问 & §9.2
\\
p61--p65 & 短路/开路规则与四张分解图 & §9.2、§9.3
\\
p66 & 1.9.3 三步 & §9.4
\\
p67 & 1.9.4 四点注意 & §9.4
\tabbot
\end{tabular}

\flowerrule

{\small\color{inkgray}
\textbf{复算说明}：本讲义报出的每一个数都由 Python 独立复算过——
p55 的 $I_1=8$、$I_2=-2$、$I_3=6\,\mathrm{A}$、$U=24\,\mathrm{V}$ 与功率收支 $320\,\mathrm{W}$；
p37 的 $U_{AB}=40-2\times14=12\,\mathrm{V}$；p38--p42 的 $U_{ab}=10+5\times6=40\,\mathrm{V}$；
p36 四个子图在 $U_S=10\,\mathrm{V}$、$R=2\,\Omega$、$I=1\,\mathrm{A}$ 下的 $8/12/-12/-8\,\mathrm{V}$；
§8.4 自拟变式的 $27/8,9/8,21/8,15/8,3/4\,\mathrm{A}$ 与功率收支 $47.25\,\mathrm{W}$；
§8.1 五阶系数矩阵在三处符号核验下的结果。

\vspace{4pt}
\textbf{标 [待确认] 的四处}：
（1）p9 的 $L=7$（桥臂接法无法从文字层复原）；
（2）p2/p3 的数值答案（原图接线不能唯一确定）；
（3）p37 的 $14\,\mathrm{A}$ 与同页其余标注不自洽（按三条并联支路解出应为 $8\,\mathrm{A}$），
本讲义按题面给的 $14\,\mathrm{A}$ 计算；
（4）p56 的拓扑读法（中横线与左右竖边相交，与 $I_4$、$I_5$ 的位置矛盾）。
这四处本讲义只备案、不猜结论。

\vspace{4pt}
\textbf{标 [教材补充] 的两处}：§8.4 的那组数值（$12\,\mathrm{V}/6\,\mathrm{V}/2\,\Omega$）是为演示求解过程自拟的，
不是课件数据；§9.4 的功率反例（两个 $1\,\mathrm{A}$ 源叠加成 $4\,\mathrm{W}$ 而非 $2\,\mathrm{W}$）是为解释
p67 第 (2) 条而补的最小例子。两处都不作为考点口径。}

\end{document}


